\documentclass[11pt]{article}
\usepackage[a4paper,margin=21mm]{geometry}
\usepackage[no-math]{fontspec}
\setmainfont{Latin Modern Roman}
\newfontfamily\CJKfont{Noto Serif CJK SC}
\XeTeXlinebreaklocale "zh"
\XeTeXlinebreakskip=0pt plus 1pt
\usepackage{amsmath,amssymb,amsthm,amscd,graphicx,color}
\usepackage{xurl}
\usepackage[unicode,hidelinks]{hyperref}
\allowdisplaybreaks
\setlength\emergencystretch{3em}
\numberwithin{equation}{section}
\usepackage{mathrsfs}
\usepackage{array}
\usepackage{tikz-cd}
\usetikzlibrary{shapes}
\usetikzlibrary[patterns]
\usetikzlibrary{decorations.pathreplacing}

\usepackage{eucal} %lettres calligraphiées

\newcommand{\Tor}{\mathrm{Tor}}

\newcommand{\dd}{ \mathrm{d} }

\newcommand{\vir}[1]{\big[#1\big]^\mathrm{vir}}
\newcommand{\fvir}[2]{[\![ #1 ]\!]^{#2}}
\newcommand{\ev}{\mathrm{ev}}
\newcommand{\pt}{\mathrm{pt}}

\newcommand{\adj}{\mathrm{Adj}}
\newcommand{\tadj}{\widetilde{\mathrm{Adj}}}


\newcommand{\m}[1]{\mfk\left[ #1 \right]}
\renewcommand{\d}[1]{\dfk\left[ #1 \right]}

\newcommand{\Jac}{\mathrm{Jac}}
\newcommand{\Pic}{\operatorname{Pic}}
\newcommand{\Alb}{\operatorname{Alb}}

\newcommand{\modulo}{ \operatorname{mod} }

\renewcommand{\Im}{\operatorname{Im}}

%Commandes raccourcies pour les caractères spéciaux

\newcommand{\abcd}{\left(\begin{smallmatrix} a & b \\ c & d \\ \end{smallmatrix}\right)}

\newcommand{\tomega}{\widetilde{\omega}}

\newcommand{\PP}{\mathbb{P}}
\newcommand{\NN}{\mathbb{N}}
\newcommand{\ZZ}{\mathbb{Z}}
\newcommand{\RR}{\mathbb{R}}
\newcommand{\CC}{\mathbb{C}}
\newcommand{\EE}{\mathbb{E}}
\newcommand{\KK}{\mathbb{K}}
\newcommand{\QQ}{\mathbb{Q}}
\newcommand{\HH}{\mathbb{H}}
\newcommand{\TT}{\mathbb{T}}
\newcommand{\FF}{\mathbb{F}}
\renewcommand{\AA}{\mathbb{A}}
\newcommand{\CP}{\mathbb{CP}}

\newcommand{\AAA}{\mathscr{A}}
\newcommand{\BBB}{\mathscr{B}}
\newcommand{\CCC}{\mathscr{C}}
\newcommand{\DDD}{\mathscr{D}}
\newcommand{\tDDD}{\widetilde{\mathscr{D}}}
\newcommand{\EEE}{\mathscr{E}}
\newcommand{\FFF}{\mathscr{F}}
\newcommand{\HHH}{\mathscr{H}}
\newcommand{\PPP}{\mathscr{P}}
\newcommand{\SSS}{\mathscr{S}}
\newcommand{\XXX}{\mathscr{X}}
\newcommand{\TTT}{\mathscr{T}}
\newcommand{\WWW}{\mathscr{W}}

\newcommand{\bfr}{\mathbf{r}}
\newcommand{\bfs}{\mathbf{s}}
\newcommand{\bfl}{\mathbf{l}}

%\newcommand{\btheta}{ {\boldsymbol\theta} }
\newcommand{\btheta}{ {\underline{\theta}} }
\newcommand{\bfH}{ {\boldsymbol\mathcal{H}} }
\newcommand{\bfX}{ {\boldsymbol X} }
\newcommand{\bfD}{ {\boldsymbol D} }
\newcommand{\bfw}{ {\boldsymbol w} }
\newcommand{\bfm}{ {\boldsymbol m} }


\newcommand{\Dfk}{\mathfrak{D}}
\newcommand{\tDfk}{\widetilde{\mathfrak{D}}}
\newcommand{\Efk}{\mathfrak{E}}
\newcommand{\Pfk}{\mathfrak{P}}
\newcommand{\tEfk}{\widetilde{\mathfrak{E}}}
\newcommand{\mfk}{\mathfrak{m}}
\newcommand{\pfk}{\mathfrak{p}}
\newcommand{\hfk}{\mathfrak{h}}
\newcommand{\lfk}{\mathfrak{l}}
\newcommand{\tfk}{\mathfrak{t}}
\newcommand{\bfk}{\mathfrak{b}}
\newcommand{\gfk}{\mathfrak{g}}
\newcommand{\sfk}{\mathfrak{s}}
\newcommand{\dfk}{\mathfrak{d}}

\newcommand{\sfb}{\mathsf{b}}
\newcommand{\sfe}{\mathsf{e}}
\newcommand{\sff}{\mathsf{f}}
\newcommand{\sfq}{\mathsf{q}}
\newcommand{\sft}{\mathsf{t}}
\newcommand{\sfs}{\mathsf{s}}
\newcommand{\sfw}{\mathsf{w}}

\newcommand{\sfA}{\mathsf{A}}
\newcommand{\tsfA}{\widetilde{\mathsf{A}}}
\newcommand{\sfB}{\mathsf{B}}
\newcommand{\sfT}{\mathsf{T}}
\newcommand{\sfW}{\mathsf{W}}
\newcommand{\sfS}{\mathsf{S}}

\renewcommand{\div}{\mathrm{div}}
\newcommand{\codeg}{\mathrm{codeg}}
\newcommand{\Area}{\mathrm{Area}}

\newcommand{\A}{\mathcal{A}}
\newcommand{\D}{\mathcal{D}}
\newcommand{\tD}{\widetilde{\mathcal{D}}}
\newcommand{\E}{\mathcal{E}}
\newcommand{\tE}{\widetilde{\mathcal{E}}}
\renewcommand{\P}{\mathcal{P}}
\renewcommand{\L}{\mathcal{L}}
\newcommand{\N}{\mathcal{N}}
\renewcommand{\S}{\mathcal{S}}
\newcommand{\M}{\mathcal{M}}
\newcommand{\R}{\mathcal{R}}
\newcommand{\G}{\mathcal{G}}
\newcommand{\J}{\mathcal{J}}
\newcommand{\X}{\mathcal{X}}
\newcommand{\T}{\mathcal{T}}
\newcommand{\V}{\mathcal{V}}
\newcommand{\C}{\mathcal{C}}
\newcommand{\K}{\mathcal{K}}
\newcommand{\B}{\mathcal{B}}
\newcommand{\Q}{\mathcal{Q}}
\renewcommand{\H}{\mathcal{H}}
\renewcommand{\O}{\mathcal{O}}
\newcommand{\F}{\mathcal{F}}
\newcommand{\Y}{\mathcal{Y}}

\newcommand{\dsum}{\displaystyle\sum}
\newcommand{\dprod}{\displaystyle\prod}
\renewcommand{\leq}{\leqslant}
\renewcommand{\geq}{\geqslant}
\renewcommand{\tilde}{\widetilde}
\renewcommand{\bar}{\overline}

\newcommand{\gen}[1]{\big\langle #1 \big\rangle}

\numberwithin{equation}{section}

\newtheorem{theo}{Theorem}[section]
\newtheorem*{theom}{Theorem}
\newtheorem{prop}[theo]{Proposition}
\newtheorem{coro}[theo]{Corollary}
\newtheorem{lem}[theo]{Lemma}
\newtheorem{conj}{Conjecture}[section]

\theoremstyle{definition}
\newtheorem{defi}[theo]{Definition}
\newtheorem{prob}[theo]{Problem}
\newtheorem{remarkk}[theo]{Remark}
\newtheorem{expl}[theo]{Example}

\newfontfamily\nativebbone{DejaVu Math TeX Gyre}
\ExplSyntaxOn
\NewDocumentCommand{\mathds}{m}{\str_if_eq:nnTF{#1}{1}{\mathord{\text{\nativebbone\char"1D7D9}}}{\PackageError{nativecompat}{Unsupported~mathds~argument}{Only~verified~digit~one~is~accepted}}}
\ExplSyntaxOff\sloppy
\makeatletter\newlabel{sec-regularity}{{6}{}{Original source locator}{}{}}\makeatother
\makeatletter\newlabel{sec-modularity}{{6.3}{}{Original source locator}{}{}}\makeatother
\makeatletter\newlabel{theo-refined-decomposition-surjective-case}{{4.12}{}{Original source locator}{}{}}\makeatother
\begin{document}
\begin{center}\Large Local correlated genus-one Gromov–Witten invariants over an elliptic curve\end{center}
\noindent\textbf{Stable ID:} 680\quad\textbf{Author:} Thomas Blomme; Francesca Carocci\par
\noindent\textbf{Work:} Correlated Gromov–Witten Invariants\par
\noindent\textbf{Source:} \url{https://sigma-journal.com/2025/046/}\par
\noindent\textbf{Version:} Official SIGMA 21 (2025) journal PDF and native TeX; acquired 2026-09-30; exact bytes pinned by SHA256\par
\noindent\textbf{Rights:} CC BY 4.0. Authors retain copyright. \url{https://creativecommons.org/licenses/by/4.0/}. Reuse and typesetting adaptation; original language and mathematical content preserved.\par
\noindent\textbf{Difficulty (prior selection preserved):} {\CJKfont H2: The author enumerates covers together with line-bundle solutions and their correlator torsor, computes the selected coefficient by isogeny arithmetic, then reconstructs all torsion coefficients using prime-power multiplicativity and a recursive relation. The arithmetic refinement and geometric covering count are both retained; the final divisor sum alone is not the problem.}\par
\medskip\noindent\textbf{Editorial boundary note (not author text):} Complete original Theorem5.3, all elliptic/log-stable-map/tangency/divisibility/correlator/group-algebra assumptions, source correlated virtual-class definitions and deformation/torsor/forgetting relations, and the full local covering-map/meromorphic-section calculation, special-correlator coefficient, monodromy/order dependence, triangular induction, multiplicativity and prime-power formulas. The final proof paragraph alone would omit nearly all novel core, so those author constructions are retained. Later refined-degeneration/floor-diagram/regularity applications are referenced but excluded. A narrowly guarded mathds resolver accepts only the verified literal digit1 indicator and renders its double-struck glyph using named font DejaVu Math TeX Gyre, an explicit portable font dependency; all other arguments are rejected. Literal author formulas and indicator conditions remain unchanged. The wrapper loads fontspec with no-math so the author legacy rm declaration and bar accent retain standard original-style Computer Modern math operator fonts, while the separately named indicator font remains explicit. A wrapper paragraph line-breaking setting relaxes a0.14pt prose overflow; author text is unchanged.\par
\noindent\textbf{External original reference sec-regularity:} Original Section6, source1712–2307, gives later floor-diagram regularity applications. Mentioned only as future use of the selected local formula.\par
\noindent\textbf{External original reference sec-modularity:} Original Section6.3, beginning1934, provides later modularity discussion. Mentioned as further details for a local-formula application, not used in its proof.\par
\noindent\textbf{External original reference theo-refined-decomposition-surjective-case:} Original refined decomposition theorem, source1230–1249. Mentioned in the local-section opening as a later combination/application, not invoked in the local-invariant proof.\par
\medskip\hrule\medskip\noindent\textbf{Author text begins below.}\par
\setcounter{section}{1}\setcounter{subsection}{0}
% Original source lines 208--286
\subsection{Setting}

In this paper, we introduce and study a refinement of the relative (or logarithmic) Gromov--Witten invariants \cite{abramovich2014stable,gross2013logarithmic,kim2010logarithmic,li2001stable} of the pair $(Y=\PP_X(\O\oplus L),D)$, where $D=D^+ + D^-=\PP_X(L)\oplus\PP_X(\O)$ and $X$ is a smooth projective variety over $\mathbb C$. We fix the following discrete data:
\begin{itemize}\itemsep=0pt
 \item $g$, $n$, $m$ are positive integers,
 \item $\bfw=(w_1,\dots,w_n)$ is a $n$-tuple of non-zero integers with $\sum w_i=0$; we moreover set \smash{$b(\bfw)=\sum_{w_i>0}w_i=-\sum_{w_i<0}w_i$};
 \item $\beta\in H_2(X,\ZZ)$ is an effective homology class.
\end{itemize}
Let $\M:=\M_{g,m}\bigl(Y|D^{\pm},\beta,\bfw\bigr)$ be the moduli space of logarithmic stable maps parametrizing
\[
f\colon\ \ (C,p_1,\dots,p_m,q_1,\dots,q_n)\to Y,
 \]
where $C$ is a stable curve of genus $g$ with image in class $\beta+b(\bfw)F\in H_2(Y,\ZZ)$, the marked point $q_i$ is mapped to $D^{\mathrm{sgn}(w_i)}$ with tangency order $|w_i|$, the marked points $p_i$ are mapped to the interior. The vector $\bfw$ is thus called \textit{tangency profile}.
The logarithmic (or equivalently relative~\cite{abramovich2014comparison}) Gromov--Witten invariants are defined integrating constraints pulled-back along evaluation maps~${
\ev \colon \M_{g,m}\bigl(Y|D^{\pm},\beta,\bfw\bigr) \longrightarrow X^n\times Y^m}$,
 against the virtual fundamental class~\smash{$\vir{\M}$}~\cite{gross2013logarithmic}.

 We propose a refinement for these GW invariants by remembering a discrete quantity depending on the position of the points mapped to the boundary divisor. This quantity is called a \textit{correlator}.

\subsection{The Albanese evaluation}
The idea of the refinement comes from reinterpreting the map to $\bigl(Y, D^\pm\bigr)=\bigl(\PP_X(\O\oplus L), D^\pm\bigr)$ as the data of a map $f\colon C\to X$ together with the choice of an isomorphism $f^*L\cong\O_C(\alpha)$ where $\alpha\in H^0\bigl(C,\bar{M}_C^{\rm gp}\bigr)$ is defined using the logarithmic structure on $C$. When $C$ is smooth, this condition takes the familiar form $f^*L\simeq\O_C(\sum w_i q_i)$.

\subsubsection[Case of X times P\^{}1]{Case of $\boldsymbol{ X\times\PP^1}$}


To simplify, let us momentarily assume that $L$ is the trivial bundle and that the source curve~$C$ is smooth.
 Then, reinterpreting $f\in\M$ as explained above, part of the data of a log stable map is the isomorphism
$
\O_C\bigl(\sum w_iq_i\bigr)\cong\O\in\Pic^0(C)$.
For a smooth curve, the Abel--Jacobi theorem ensures that $\Pic^0(C)\simeq\Alb(C)$, where $\Alb(C)$ is the Albanese variety of $C$.
 We denote by~${a_X\colon X\to\Alb(X)}$ the map from $X$ to its Albanese variety, unique up to translation. By the functoriality of the Albanese variety, $f$ induces a~morphism of Abelian varieties
$
 f_*\colon \Alb(C)\to\Alb(X)
$ compatible with the Albanese maps.

We denote by $a_\bfw$ the morphism
\[
 a_\bfw\colon\ (x_1,\dots,x_n)\in X^n\longmapsto \sum_{i=1}^n w_i a_X(x_i)\in\Alb(X).
 \]
As $\sum w_i=0$, the map $a_\bfw$ does not change if we compose $a_X$ with a translation. From $\O_C(\sum w_iq_i)\simeq\O\in\Pic^0(C)$, we get that
\begin{equation}\label{eq-f*f*}
 \sum_{i=1}^n w_i a_X(f(q_i)) =0\in\Alb(X).
\end{equation}
Equation \eqref{eq-f*f*} is a constraint on the relative position of the images $f(q_i)$. In other words, the~$X^n$ part of the evaluation map is not surjective, and truly has values in the subset $a_\bfw^{-1}(0)$.


\subsubsection{Case of a non-trivial bundle} If $L$ is not assumed to be trivial anymore, we prove in Lemma~\ref{lem:mapphibeta} that there is a natural map $\varphi_\beta\colon \Pic^0(X)\to\Alb(X)$ defined through Hodge theory, which only depend on the curve class. Then
Proposition \ref{prop-image-of-L-is-constant} ensures that equation \eqref{eq-f*f*} becomes the following:
\begin{equation}\label{eq-f*f*L}
\sum_{i=1}^n w_ia_X(f(q_i)) = \varphi_\beta(L)\in\Alb(X).
\end{equation}
In other words, the image depends only on the degree $\beta$ and the line bundle $L$ and not on $f$.

\subsection{Correlated classes and correlated GW invariants} We now assume that the tangency orders $w_i$ have a non-trivial common divisor $\delta$. We then consider the map
\[
a_{\bfw/\delta}\colon\ (x_1,\dots,x_n)\in X^n\longmapsto\sum_1^n\frac{w_i}{\delta}a_X(x_i)\in \Alb(X).
 \]
The above morphism induces a morphism from the space of log stable maps $\M$
\[\kappa^\delta(f)=\sum \frac{w_i}{\delta}a_X(f(q_i))\in\Alb(X).\]
It follows from equation \eqref{eq-f*f*L} that
$\delta\cdot\kappa^\delta(f)=\varphi_\beta(L)$.
Therefore, the map $\kappa^\delta$ takes values in the set of $\delta$-roots of $\varphi_\beta(L)$, denoted by $T^\delta(L,\beta)$. The latter is a torsor under the group of $\delta$-torsion elements in $\Alb(X)$, denoted by $\Tor_\delta(\Alb(X))$.
The value of $\kappa^\delta$ is called a \textit{correlator} and is denoted by $\theta$.
We want to refine the Gromov--Witten of $\bigl(Y,D^{\pm}\bigr)$ keeping track of the value of the correlator.

The discussion above assumes the source curve $C$ to be smooth, but the morphism $\kappa^\delta$
is also defined on the boundary of the moduli space $\M_{g,m}\bigl(Y|D^{\pm},\beta,\bfw\bigr)$, still with values in $T^\delta(L,\beta)$ (see Section~\ref{sec:correlationrefinement}).
This shows that
 $\M$ splits into distinct (possibly still disconnected) components~$\M^\theta$ indexed by the values of the correlator $\theta\in T^\delta(L,\beta)$; in particular the virtual class splits as a sum of the so-called \textit{correlated classes} \smash{$\vir{\M^\theta}$}. We define the \textit{full correlated class}~$\fvir{\M}{\delta}$
\[ \fvir{\M}{\delta}=\sum_{\theta\in T^\delta(L,\beta)} \vir{\M^\theta}\cdot(\theta)\in \QQ[\Alb(X)]\otimes H_\bullet(\M,\QQ), \]
which is an element of the group algebra with cycle coefficients with support on $T^\delta(L,\beta)\subset\Alb(X)$.

\textit{Correlated GW invariants} are obtained by capping these classes with some pullback by the evaluation map: if $\gamma\in H^\bullet(X^n\times Y^m,\QQ)$, we set
\[ \langle\langle\gamma\rangle\rangle^\delta_{g,\beta,\bfw} = \int_{\fvir{\M}{\delta}}\ev^*\gamma \in\QQ[\Alb(X)], \]
which is an element of the group algebra $\QQ[\Alb(X)]$ with support on the torsor $T^\delta(L,\beta)$. Distinct correlators may yield different Gromov--Witten invariants, as illustrated by the computations in the elliptic case.

It follows from the discussion above that the evaluation map truly has values in the set \smash{$a_\bfw^{-1}(\varphi_\beta(L))=\bigsqcup a_{\bfw/\delta}^{-1}(\theta)$}. A broader generalization of GW invariants would be obtained by allowing the pullback of cohomology classes from $a_\bfw^{-1}(\varphi_\beta(L))$ rather than $X^n$. Correlated invariants are a particular case of these invariants where we pull-back the classes $1_\theta\in H^0\bigl(a_\bfw^{-1}(\varphi_\beta(L)),\QQ\bigr)$, corresponding to the components $a_{\bfw/\delta}^{-1}(\theta)$ of $a_{\bfw}^{-1}(\varphi_\beta(L))$.




% Original source lines 369--959
\section{Logarithmic curves, line bundles and stable maps}
\subsection{Curves and line bundles}
We
refer the reader to \cite{ogus} for an extensive introduction to logarithmic structures and in particular for all the basic definitions which we do not recall in what follows.
\subsubsection{Logarithmic curves}
Let $(S,M_S)$ be a fine and saturated logarithmic scheme. A log curve over $S$ is a proper, integral and logarithmically smooth morphism $\pi\colon C\to S$ with connected, reduced one-dimensional (geometric) fibers.
Kato provided the following local characterization.
\begin{theo}[\cite{KatoF}]
For every geometric point $p\in C$ with $\pi(p)=s$, there exists an \'etale local neighbourhood of $p$ in $C$ with a strict \'etale morphism to
\begin{description}\itemsep=0pt
 \item[\emph{smooth point:}] $\mathbb A^1_S$ with log structure pulled back form the base, i.e., $\overline{M}_{C,p}=\overline M_{S,s}$;
 \item[\emph{marking:}] $\mathbb A^1_S$ with log structure generated by the $0_S$-section and $\pi^*M_S$, i.e., $\overline{M}_{C,p}=\overline M_{S,s}\oplus \mathbb N$;
 \item[\emph{node:}] $\operatorname{Spec}(\mathcal O_S[x,y]/xy-t)$ for some $t\in\mathcal O_S$, with log structure induced by the multiplication map $\mathbb A^2_S\to\mathbb A^1_S$ and $t\colon S\to \mathbb A^1$. In particular, $\overline{M}_{C,p}=\overline M_{S,s}\oplus \mathbb N e_x\oplus \mathbb N e_y\slash e_x+e_y=\delta$ for~${\delta\in \overline{M}_{S,s}}$ the so called \emph{smoothing parameter}.
\end{description}
\end{theo}

\subsubsection{Tropical curves}

To a family of logarithmic curves, we can associate a family of tropical curves.
Consider first $C\to (\operatorname{Spec}(k),Q)$ a logarithmic curve over a log point.

The \emph{tropicalization} $\Gamma$ is defined as the generalised cone complex obtained as the following colimit:
\[
\varinjlim_{\eta\rightsquigarrow\xi} \operatorname{Hom}\bigl(\overline{M}_{C,p},\mathbb R_{\geq 0}\bigr),
\]
where $\eta\rightsquigarrow\xi$ denote specialization maps inducing surjective morphisms $\overline{M}_{C,\xi}\to\overline{M}_{C,\eta}$. From the local description of the log structure on log smooth curves recalled above, it follows that $\Gamma$ comes equipped with a morphism of cone complexes
$\Gamma\to\sigma_Q $,
for $\sigma_Q=\operatorname{Hom}(Q,\RR_{\geqslant 0})$ the cone dual to the monoid $Q$.

Alternatively, $\Gamma$ can be thought as the dual graph of $C$ (having a vertex for each irreducible component, a leg for each marking and an edge $e\in\Gamma$ for each node) together with a metrization of the bounded edges with values in the base monoid $Q$. This means that we have an assignment~${l\colon E_b(\Gamma)\to Q}$ defined by associating to each bounded edge its smoothing parameter~${\delta_e\in Q}$.

\begin{remarkk}
Notice that for each map $Q\to\mathbb R_{\geq 0}$ we obtain a tropical curve in the classical sense. Since $Q\to\mathbb R_{\geq 0}$ correspond to a point of the cone $\sigma_Q$ dual to $Q$ (in the usual sense of toric geometry), we can think of a tropical curve metrized in $Q$ as a family of usual tropical curves on the dual cone.
\end{remarkk}
\begin{remarkk}\label{rem:tropfamilies}
The tropicalization makes sense over more general base log schemes (see also \cite[Definition~2.3.3.3]{{molchowiselog}}). Let $(C,M_C)\to(S,M_S)$ a log smooth curve;
for each $s\in S$, $\Gamma_s$ consists of the dual graph of $C_s$ metrized in $\overline{M}_{S,s}$.
 For each specialization $\eta\rightsquigarrow s$, we have a compatible diagram
\[
\begin{tikzcd}
\Gamma_s\ar[r,"c"]\ar[d,"\ell"] &\Gamma_{\eta}\ar[d,"\ell"]\\
\overline{M}_{S,s}\ar[r] & \overline{M}_{S,\eta},
\end{tikzcd}
\]
where $\overline{M}_{S,s}\to \overline{M}_{S,\eta}$ is a localization morphism followed by sharpification (quotient the invertibles) and $c$ is an edge contraction; more precisely, $c$ contracts the edges of $\Gamma_s$ whose length becomes zero in $ \overline{M}_{S,\eta}$.

\end{remarkk}
\subsubsection{Piecewise linear functions}
Unravelling the definition of global sections of the characteristic sheaf, we get a convenient description of $H^0\bigl(C,\overline{M}_C^{\rm gp}\bigr)$ in terms of continuous functions on the tropicalization $\Gamma$ which are linear with integral slope along the edges.
Let us first assume that $C$ is a log curve over a~${S=(\operatorname{Spec}(k),Q)}$.

Since $\overline{M}_C^{\rm gp}$ is a constructible sheaf, a
section is determined by giving $m_x\in\overline{M}_{C,x}^{\rm gp}$ at a generic point $x$ of each stratum,
such that the values $m_x$ are compatible with specialization, i.e., if $\xi$ generalize to $\eta$ then the image of $m_{\xi}$ under the surjection $\overline{M}_{C,\xi}^{\rm gp}\to \overline{M}_{C,\eta}^{\rm gp}$ is $m_{\eta}$.

Looking at Kato's classification, there are three type of strata on log smooth curves: the open strata corresponding to irreducible components of the curves excluding markings, i.e., vertices in the tropicalization; the closed strata corresponding to markings; the closed strata corresponding to nodes. Therefore, a section $\alpha\in H^0\bigl(C,\overline{M}_C^{\rm gp}\bigr)$ is the data of
\begin{itemize}\itemsep=0pt
 \item for each vertex $V$ of $\Gamma$ a value $\alpha(V)\in Q$;
 \item for each leg and edge an integer $w_e\in\mathbb Z$ (the slope along edge or leg in $\Gamma$);
 \item a compatibility condition for the values $\alpha(V)$, $\alpha(W)$ of vertices incident to the same edge~$e$: $\alpha(V)-\alpha(W)=\delta_e w_e$.
\end{itemize}
Indeed, suppose that $x_e$ is a node
of $C$ which generalizes to the generic points $\eta_V$ and $\eta_W$; then we have a map
$\overline{M}_{C,x_e}^{\rm gp}\to Q\times Q$
that is injective and gives an identification
\[\overline{M}_{C,x_e}^{\rm gp}=\{(a,b)\in Q\times Q \mid a-b\in\delta_e\mathbb Z\}.\]

We denote by $\operatorname{CL}(\Gamma,Q)$ the set of functions satisfying the above conditions, called \textit{cone-wise linear}.\footnote{We reserve the name piece-wise linear for those function which become linear along the edges after a subdivision.} The previous discussion can be summarized in the following equality (see also \cite[Section~3.3.3]{marcuswiselog} and references therein)
$H^0\bigl(C,\overline{M}_C^{\rm gp}\bigr)=\operatorname{CL}(\Gamma,Q)$.

 \begin{remarkk}
 Using the description of tropicalization for a family $(C,M_C)\xrightarrow{\pi}(S,M_S)$ of log curves given in Remark~\ref{rem:tropfamilies} and the above discussion, we can think of a section $ \alpha\in H^0\bigl(S,\pi_*\overline{M}_C^{\rm gp}\bigr)$ as a collection \smash{$\alpha_s\in \operatorname{CL}\bigl(\Gamma_s,\overline{M}_{S,s}^{\rm gp}\bigr)$} compatible under edge contraction, i.e., generalization.
 \end{remarkk}



\subsubsection{Line bundles}
Let $C\xrightarrow{\pi} S$ be a log smooth curve. We have a short exact sequence
\begin{equation}\label{eq:logseq}
0\to\mathcal O_C^{\times}\to M_C^{\rm gp}\to\overline{M}_C^{\rm gp}\to 0,\end{equation}
from which we obtain a long exact sequence of sheaves on $S$
\begin{align*}
 \cdots\to \pi_*M_C^{\rm gp}\to\pi_*\overline{M}_C^{\rm gp}\to
 R^1\pi_*\mathcal O_C^{\times}\to
 R^1\pi_*M_C^{\rm gp}\to R^1\pi_*\overline{M}_C^{\rm gp}\to\cdots.
\end{align*}
In particular, if the underlying scheme of $S$ is a point, so that $\pi_*\overline{M}^{\rm gp}_C=H^0\bigl(C,\overline{M}^{\rm gp}_C\bigr)$ and $R^1\pi_*\O_C^\times=H^1\bigl(C,\O_C^\times\bigr)=\Pic(C)$, we get a map
\smash{$
 H^0\bigl(C,\overline{M}_C^{\rm gp}\bigr)\to \Pic^0(C)$}, $
 \alpha \mapsto\mathcal O_C(-\alpha)$.

\begin{remarkk}
 The sign comes from observing that when $\alpha\in H^0\bigl(C,\overline{M}_C\bigr)$, the $\mathcal O_C^{\times}$-torsors of lifts of $\alpha$ to a section of $M_C$ comes equipped with a map to $\mathcal O_C$, coming from the structure morphism~${\epsilon \colon M_C\to\mathcal O_C}$ of the log structure. Up to changing $\alpha$ by $-\alpha$, we now deal with $\O_C(\alpha)$ instead.
\end{remarkk}

The restriction of the line bundle $\mathcal O_C(\alpha)$ to the irreducible components $C_V$ of $C$ can be explicitly described \cite[Proposition~3.3.3]{marcuswiselog}.
\begin{prop}
 Let $C\xrightarrow{\pi} S=(\operatorname{Spec}(k),Q)$ be a log curve and $\alpha\in\operatorname{CL}(\Gamma, Q)$, where $\Gamma$ denotes the tropicalization of $C$. Then for any vertex $V$,
 \[\mathcal O_C(\alpha)\rvert_{C_V}=\pi^*\mathcal O_S(\alpha(V))\otimes\mathcal O_{C_V}\biggl(\sum_{e\vdash V} s_e(\alpha) q_e\biggr),
 \]
where the sum is taken over all the edges and legs $e$ incident to $V$ and $s_e(\alpha)$ is the slope of~$\alpha$ along~$e$, oriented away from~$V$, where $q_e$ denotes the preimage of the node associated to the edge~$e$ on the considered component.
\end{prop}
\begin{remarkk}
 The result is more generally true if $C\to S$ is a logarithmic curve with constant degeneration.
\end{remarkk}

\subsection{Logarithmic line bundles and trivializations}
\begin{defi}[\cite{molchowiselog}]
 Let $C\to S$ be a log smooth curve. A \emph{logarithmic line bundle} $P$ on $C$ is a $M_C^{\rm gp}$-torsor such that its restriction $P\rvert_{C_s}$ to the geometric fibers has \emph{bounded monodromy}.
\end{defi}
We refer the reader to \cite{molchowiselog} for all explanations about the meaning and the necessity of the bounded monodromy. Here we only need that it is a condition on the $\overline{M}_C^{\rm gp}$-torsor $\overline{P}$ associated to $P$, and that it is automatically satisfied when $\overline{P}$ is isomorphic to the trivial $\overline{M}_C^{\rm gp}$-torsor.

In particular, we see from the long exact sequence
\[H^0\bigl(C,\overline{M}_C^{\rm gp}\bigr)\to\Pic(C)\to H^1\bigl(C,M_C^{\rm gp}\bigr)\to H^1\bigl(C,\overline{M}_C^{\rm gp}\bigr)\to\cdots\]
that for any $L\in \Pic(C)$ the associated $M_C^{\rm gp}$-torsor \smash{$L^{\times}\otimes_{\mathcal O_C^{\times}}M_C^{\rm gp}$} is in fact a logarithmic line bundle, where we denote by $L^{\times}$ the $\mathbb{G}_m$-torsor obtained from the line bundle removing its zero section.

\begin{defi}[{\cite[Definition~4.5.1]{marcuswiselog}}]
 Let $C\xrightarrow{\pi} S$ be a log smooth curve and $L$ a line bundle on $C$. A \emph{logarithmic trivialization} of $L$ is a section of \smash{$L^{\times}\otimes_{\mathcal O_C^{\times}}M_C^{\rm gp}$}.
\end{defi}
The definition is a modification of \cite[Definition~4.5.1]{marcuswiselog}, where the authors define logarithmic trivializations relative to the base. The difference comes from the fact that \cite{marcuswiselog} are interested in moduli of logarithmic maps to rubber targets, while we study the rigid case.
Similarly, the following discussion is parallel to
\cite[Proposition~4.5.3]{marcuswiselog}.

From \eqref{eq:logseq}, we see that the logarithmic line bundle \smash{$L^{\times}\otimes_{\mathcal O_C^{\times}}M_C^{\rm gp}$} can be thought as a \smash{$\mathcal O_C^{\times}$}-torsor over the trivial $\overline{M}_C^{\rm gp}$-torsor.
In particular, there is an exact sequence
\[0\to H^0\bigl(C,L^{\times}\otimes_{\mathcal O_C^{\times}}M_C^{\rm gp}\bigr)\to H^0\bigl(C, \overline{M}_C^{\rm gp}\bigr)\to\Pic(C),
 \]
where the second map is defined by
$\alpha\mapsto L(-\alpha)=L\otimes\O(-\alpha)$.
In other words, a logarithmic trivialization of $L$ is a trivialization of the line bundle $L(-\alpha)$ for some cone-wise linear function~$\alpha$.

\subsection[Logarithmic stable maps to P\^{}1-bundles and their degenerations]{Logarithmic stable maps to $\boldsymbol{\mathbb P^1}$-bundles and their degenerations}

This subsection follows closely \cite[Section~5]{marcuswiselog}, with two small differences. First, \cite{marcuswiselog} considers maps to rubber target, while we deal with the rigid case. Second, they study maps to $\mathbb{P}^1$ or to degeneration of $\mathbb P^1$ to a chain, while we consider maps to $\mathbb{P}^1$-bundles $Y=\mathbb P_X(\mathcal O_X\oplus L)$ and degenerations of the latter to certain chains $\Y_0=\bigl(Y_1,D_1^\pm\bigr)\cup \bigl(Y_2,D_2^\pm\bigr)\cup\dots\cup\bigl(Y_N,D_N^\pm\bigr)$ of~$\mathbb{P}^1$-bundles over $X$ obtained by gluing $Y_i$ to $Y_{i+1}$ along $D_i^+\cong D_{i+1}^-\cong X$.

For the benefit of the reader, we state the results in the form most convenient for us and sketch how to adapt the proofs of \cite[Section~5]{marcuswiselog} to our situation.

Let $X$ be a smooth projective variety and $L$ a line bundle on $X$.
For a fixed $\beta\in H_2(X,\mathbb Z)$ and a vector of integers $\bfw=(w_1,\dots,w_n)$ such that $\sum_{i=1}^n w_i=c_1(L)\cdot\beta$, we consider the moduli space
$\M := \M_{g,m}\bigl(Y|D^{\pm},\beta,\bfw\bigr)$
of logarithmic stable maps to $Y=\mathbb P_X(\mathcal O_X\oplus L)$ endowed with the divisorial log structure $D^++D^-=\mathbb P_X(\mathcal O_X) +\mathbb P_X(L)$ and with prescribed contact order $w_i$ at $q_i$ along $D^++D^-$ \cite{abramovich2014stable,gross2013logarithmic}.
Positive $w_i$ encode the tangencies with $D^+$ and the negative $w_i$ the tangencies with $D^-$.

This space compactifies the moduli space parametrizing maps from marked curves to $X$,
\[f\colon\ \bigl(C,p_1,\dots, p_m,q_1,\dots,q_n\bigr)\to X,\]
together with an isomorphism
\[\lambda\colon\ f^* L\biggl(-\sum a_i q_i\biggr)\cong\mathcal O_C.\]

Using the language of logarithmic trivializations introduced in the previous subsection, this description can be extended over the boundary.

\begin{defi}\label{def:logstablemaps}
 We define $\M(X,L)$ to be the stack over $\operatorname{LogSch}^{fs}$ whose objects are
 \begin{enumerate}\itemsep=0pt
 \item [(1)] A log smooth curve $C\to S$ with $n$ marked points with logarithmic structure and $m$ schematic marked points.
 \item [(2)] A stable map $f\colon (C,p_i,q_i)\to X$.
 \item [(3)] A logarithmic trivialization \smash{$\lambda\in H^0\bigl(C,f^*L^{\times}\otimes_{\mathcal O_C^{\times}}M_C^{\rm gp}\bigr)$}, such that the induced $\alpha\in H^0\bigl(C,\allowbreak\overline{M}_C^{\rm gp}\bigr)$ is locally on $C$ comparable with $0$ and the slope of $\alpha$ along the $i$-th marking is $w_i$.
 \end{enumerate}

We recall that a section \smash{$\alpha\in H^0\bigl(C,\overline{M}^{\rm gp}_C\bigr)$} is said comparable to $0$ if at each point $x\in C$ we have that either $\alpha(x)\geq 0$ or $\alpha(x)\leq 0$ in the partial order on $\overline{M}^{\rm gp}_{C,x}$ defined by the monoid $\overline{M}_{C,x}$. We refer the reader to \cite[Section~2.5]{RSPW2} for further explanations on the geometric meaning of this condition.
\end{defi}

\begin{prop}\label{prop:logstabelmaps}
The moduli space $\M(X,L)$ is equivalent as a stack over logarithmic schemes to the moduli space of logarithmic stable maps $\M$ of Gross--Siebert, Abramovich--Chen with target~${(Y=\PP_X(\mathcal O_X\oplus L),\mathcal M_Y)}$ and fixed tangencies $w_i$ along the markings $q_1,\dots,q_n$ {\rm\cite{abramovich2014stable,gross2013logarithmic}}. Here~$\mathcal M_Y$ is the divisorial log structure defined by the fiberwise $0$ and $\infty$ sections.
\end{prop}
\begin{proof}
The proof essentially consists in unravelling the definition of logarithmic stable map to a $\mathbb P^1$-bundle. For the case of log stable maps to toric varieties, the analogous result is proved, for example, in \cite[Proposition~2.5.1.1]{RSPW2}.

Let $\mathrm{Tot}(L)\xrightarrow{\pi} X$ be the total space of $L$ endowed with the zero section logarithmic structure~$M_L$. First, we argue that a logarithmic stable map
\smash{$C\xrightarrow{F} \mathrm{Tot}(L)$} is the same as a logarithmic trivialization $\lambda$ such that the associated $\alpha$ is a section of the log structure, i.e., $\alpha\in H^0\bigl(C,\bar{M}_C\bigr)$.
Indeed, by definition, a log map is a map $\underline{F}$ of the underlying schemes, namely $f\colon C\to X$ plus a section of $f^*L$, together with a morphism $\underline{F}^{-1}M_L\to M_C$ of the logarithmic structures. Using the Borne--Vistoli description of fine and saturated log structures \cite{BorneVistoli}, the latter is the same as a~morphism of the characteristic sheaves $\underline{F}^{-1}\bar{M}_L\to \bar{M}_C$ such that the maps to $\mathfrak{D}\mathrm{iv}_C$ commute. The first map comes from
$
 \bar{M}_L\cong\mathbb N_X\to \mathfrak{D}\mathrm{iv}_{\mathrm{Tot}(L)}$, $
1\to \mathcal O_L(X)\cong \pi^* L
$
pulling back along $\underline{F}$. It follows that a lift of the morphism of characteristic sheaves to a~morphism of the logarithmic structures correspond to an isomorphism $f^*L\cong^{\lambda}\mathcal O_C(\alpha)$ for $\alpha$ the section of $\bar{M}_C$ induced by the map on characteristic sheaves.

At this point we argue parallel to \cite[Section~2.5.1]{RSPW2}.

The $\mathbb P^1$-bundle $\mathbb P_X(\mathcal O_X\oplus L)$ with its fiberwise toric log structure can be constructed as the quotient of the rank 2 vector bundle $\mathcal O_X\oplus L$ minus the $0$ section for the fiberwise $\mathbb G_m$-action.

Then any logarithmic map $C\to \mathbb P_X(\mathcal O_X\oplus L)$ locally lifts to $\O_X\oplus L\setminus 0_X$.

By the discussion above, this lift can be represented by $a$, $b$ logarithmic trivializations of~$\mathcal O_C$ and $f^*L$ such that the associated $\bar{M}_C^{\rm gp}$ section $\bar{a}$ and $\bar{b}$ are in $\bar{M}_C$. Notice that the ratio~${\lambda=ab^{-1}}$ is a well defined $\mathbb G_m$-invariant section of the torsor \smash{$f^*L^{\times}\otimes_{\mathcal O_C^{\times}}M_C^{\rm gp}$}, namely a logarithmic trivialization of $f^*L$.

Notice however that not all the section $\alpha\in \rm{H}^0\bigl(C,\bar{M}_{C}^{\rm gp}\bigr)$ can arise this way: indeed, since $(a,b)$ must avoid the zero section, this means that locally around each point $x\in C$ either $\bar{a}_x\in \bar{M}_{C,x} $ or $\bar{b}_x\in\bar{M}_{C,x}$ is zero (i.e., lifts to an isomorphism of trivial $\mathbb G_m$-torsors). For $\alpha=\bar{a}-\bar{b}$, this translates precisely to being locally comparable to zero.
\end{proof}

\begin{remarkk}
If in Definition~\ref{def:logstablemaps}, we drop the request on the local comparability with zero of the section $\alpha\in H^0\bigl(C,\bar{M}_C^{\rm gp}\bigr)$, we would obtain the stack over $\operatorname{LogSch}^{fs}$ parametrizing log maps to the $\mathbb G_{log}$ bundle $\rm{Tot}(L)\times_{\mathbb G_m}\mathbb G_{log}$ associated with $L$. Similar moduli stack have been considered various times in the literature \cite{marcuswiselog,ranganathan2024logarithmic,RSPW2,ranganathan2020rational}.\footnote{We thank an anonymous referee whose comments suggested us to add this remark.}
\end{remarkk}

\begin{remarkk}
The definition and the characterization above can easily be adapted to the case of moduli of logarithmic stable maps to expansions \cite{kim2010logarithmic,li2001stable,maulik2023logarithmic}. Indeed, it is sufficient to add the condition that the values of $\alpha(V)$ are totally ordered; similar discussions appear in \cite{marcuswiselog,ranganathan2024logarithmic}.

In the classical language of \cite{li2001stable}, maps in the boundary of $\M$ which differ by the rubber action are identified. It is explained in detail in \cite{CN} how to see this from the logarithmic prospective.
\end{remarkk}

We are also interested in considering log maps to degenerations, i.e.,
\begin{itemize}\itemsep=0pt
\item Logarithmic maps to \smash{$\Y_t\xrightarrow{p}\mathbb A^1_t$} certain one parameter degenerations of $\mathbb P_{\X_{t\neq 0}}(\O\oplus\L)$ to
\[\Y_0=\bigl(Y_1,D_1^\pm\bigr)\cup\bigl(Y_2,D_2^\pm\bigr)\cup\dots\cup \bigl(Y_N,D_N^\pm\bigr),
\]
where $\bigl(Y_i,D^{\pm}_i\bigr)$ are $\mathbb P^1$-bundles over $\X_0$ glued along their infinity and zero sections;
$\mathbb A^1_t$ is endowed with its toric log structure and $\Y_t$ with the divisorial log structure defined by~${\Y_0+\mathcal{D}^+ +\mathcal{D}^-}$, where $\mathcal{D}^\pm$ are the $0$ and $\infty$-section of the family, and making $p$ log smooth.
\item Logarithmic maps to the log smooth scheme $\Y_0\to(\operatorname{Spec} k,\mathbb N)$ with the log structure pulled back from the family $\Y_t$
\end{itemize}

We recall the following definition.

\begin{defi}[{\cite[Definition~5.2.1]{marcuswiselog}}]
 A \emph{divided tropical line} over a logarithmic scheme $(S,M_S)$ is defined as follows. Let $P$ be an $M_S^{\rm gp}$-torsor and $\bar{P}$ the associated $\bar{M}_S^{\rm gp}$ torsors; fix $\gamma_0\leq\gamma_1\leq\cdots\leq\gamma_N$ a non empty collection of sections of $\bar{P}$ defined locally on $S$. Then $\bar{P}_{\gamma}\subset\bar{P}$ is defined to be the subfunctor whose local sections are comparable with all the $\gamma_i$.
\end{defi}
It is proved in \cite[Proposition~5.2.4]{marcuswiselog} that $P_{\gamma}:=\bar{P}_{\gamma}\times_{\bar{P}} P$ is a 2-marked family of semi-stable genus zero curve.



We then consider the following variation: fix $(\mathcal X, M_{\X})\xrightarrow{p} (S,M_S)$ with $X\to S$ smooth projective and $M_{\X}=p^*M_S$.
The two cases of interest listed above correspond to the case where~$(S,M_S)$ is $\mathbb A^1$ with its toric log structure or a standard log point.

Let $P\in\operatorname{LogPic}_{\X\slash S}$ and
$\bar{P}$ the induced $\bar{M}^{\rm gp}_{\mathcal{X}}$-torsor. Let $\gamma_0\leq\gamma_1\leq\cdots\leq\gamma_N$ be a non empty collection of sections of \smash{$\bar{P}$}. Notice that where the logarithmic structure of $\X$ is trivial $P$ is a line bundle in the usual sense and all the $\gamma_i$ coincide and are necessarily zero.
\begin{lem}\label{lem:tropicaline}
In the notation above, let $\bar{P}_{\gamma}\subseteq \bar{P}$ be the sub-funtcor of sections locally comparable with all the $\gamma_i$ and $P_{\gamma}=P\times_{\bar{P}}\bar{P}_{\gamma}$. Then $P_{\gamma}\to\X\to S$ is an expanded degeneration of $\Y=\mathbb P_{\X}(\L\oplus\O)$ for $\L$ the line bundle on $\X$ of lifts of $\gamma_0$ to a section of $P$.
\end{lem}
\begin{proof}
Since we assumed that the collection of sections is not empty, $\bar{P}$ admits a trivialization~${\gamma_0\colon\X\to\bar{P}}$ which we think as the zero of $\bar{M}_{\X}$ group. Notice that if the log structure on~$S$ and thus on $\X$ is generically trivial, then all the $\gamma_i$ coincide and are actually zero on the locus of $\X$ with trivial logarithmic structure. Then there exist a line bundle $\L_0$ representing the log line bundle $P$, namely the $\O_{\X}^\times$-torsor of lifts of $\gamma_0$ to a section of the log line bundle. We already proved in Proposition~\ref{prop:logstabelmaps} that the subfunctor of \smash{$P=L_0^\times\otimes_{\O_{\X}^\times}M_{\X}^{\rm gp}$} whose local sections are~comparable with $\gamma_0$ is the $\mathbb{P}^1$-bundle $\Y=\mathbb P_{\X}(\L_0\oplus\O)$ endowed with its boundary logarithmic structure.
Once a trivialization is fixed, we have that
\[
\gamma_i-\gamma_0=\delta_i\in H^0\bigl(\bar{M}_{\X}^{\rm gp}\bigr);
\] these are maps \smash{$\operatorname{Trop}(\X)=\operatorname{Trop}(S)\xrightarrow{\delta_i}\operatorname{Trop}(\Y)$}. Subdividing along the image, we obtained the desired degeneration.

 Alternatively, the statement can be proved following the steps of \cite[Proposition~5.2.4]{marcuswiselog}.
\end{proof}

Let $\mathcal Y_t\to\mathbb A^1$ a semi-stable degeneration coming form a subdivided tropical line; for example, this is the case if $\Y_t$ is constructed starting from $\mathbb P_{\X_t}(\O_{\X_t}\oplus\L_t)$ by successively blowing up the zero and the infinity section on the central fiber.
Denote by $\M\bigl(\Y_t\slash\mathbb A^1\bigr)$ the moduli space of logarithmic stable maps to $\Y_t\slash\mathbb A^1$ from families of log smooth curves with fixed genus $g$, number of markings $m$, contact order $(w_1,\dots, w_n)$ along $\D_t^-+\D_t^+$ and curve class $\beta$. Then we have the following.

\begin{coro}\label{prop:logmapsdegeneration}
 The moduli space $\M\bigl(\Y_t\slash\mathbb A^1\bigr)$ parametrizes the following data:
 \begin{enumerate}\itemsep=0pt
 \item[$(1)$] A diagram of logarithmic maps
 \[
 \begin{tikzcd}
 C\ar[r,"f"]\ar[d] & \X_t\ar[d]\\
 S\ar[r] &\mathbb A^1
 \end{tikzcd}
 \]
 for $C\to S$ a family of log smooth curves.
 \item[$(2)$] A section \smash{$\lambda\in H^0\bigl(C, f^* \mathcal L^{\times}\otimes_{\mathcal O^{\times}_{C}}M_{C}^{{\rm gp}}\bigr)$} such that the induces $\alpha\in H^0\bigl(C,\bar{M}_C^{\rm gp}\bigr)$ is locally comparable with $f^*\gamma_i$.
 \end{enumerate}
\end{coro}
\begin{remarkk}
 Parallel to before, if we want to instead consider logarithmic maps to expanded degenerations \cite{kim2010logarithmic,li2001stable,maulik2023logarithmic}, it suffices to further impose that for each geometric point $s$ the values of $\alpha(V)$ for $V\in\Gamma_s$ are totally ordered.
\end{remarkk}

\section{Correlated virtual class}

\subsection{Recollection on Albanese varieties}

\subsubsection{Complex setting}
For $X$ a proper, smooth variety over the complex numbers the Albanese variety was originally defined via transcendental methods using path integrals of closed holomorphic 1-forms. Let~$H_1(X,\mathbb Z)$ denote the torsion free part of the first homology group. There is an inclusion via the integration on paths:
\begin{align*}
 H_1(X,\mathbb Z)\longrightarrow H^0(X,\Omega_X)^*,\qquad
 \gamma\longmapsto \left(\omega\mapsto\int_{\gamma}\omega\right),
\end{align*}
where $ H^0(X,\Omega_X)^*$ denotes the dual to the vector space of holomorphic $1$-forms.
The Albanese variety of $X$ is the complex abelian variety obtained as the following quotient:
\[\Alb(X):=H^0(X,\Omega_X)^*\slash H_1(X,\mathbb Z). \]
We list below some results and properties of the Albanese variety and refer the reader to \cite[Section~11.11]{birkenhake} or \cite{mumford1974abelian} and references therein for all details and proofs.
\begin{enumerate}\itemsep=0pt
 \item[(1)] When $X$ is an Abelian variety, then $\Alb(X)\cong X$.
 \item[(2)] Let $x_0\in X$ a point, then there is an algebraic map, called the \emph{Albanese map}
 \begin{align*}
 a_{X}\colon \ X\longrightarrow \Alb(X),\qquad
 x\longmapsto\left(\omega\mapsto\int_{x_0}^x\omega\right) \mod H_1(X,\mathbb Z),
 \end{align*}
 sending $x_0$ to the identity element. This map is unique up to translation in the sense that choosing a different $x_0$ amounts to compose $a_X$ with a translation.
 \item[(3)] The Albanese map satisfies the following universal property:
 let $\varphi\colon X\to A$ be a map to an abelian variety, then there exists a unique homomorphism of abelian varieties $\widetilde{\varphi}\colon \Alb(X)\to A $ such that $\widetilde{\varphi}(0)=\varphi(x_0)$ and $\widetilde{\varphi}\circ a_{X}=\varphi$.
 \item[(4)] Let $X\xrightarrow{f} Y$ a morphism of algebraic varieties, then the map to the Albanese varieties are functorial, i.e., there exists a homomorphism of abelian varieties $f_*$ making the following diagram commute:
 \[
 \begin{tikzcd}
 X \arrow[r,"f"] \arrow[d,"a_X"']& Y \arrow[d,"a_Y"] \\
 \Alb(X) \arrow[r,"f_*"] & \Alb(Y),
 \end{tikzcd}
\]
 where $a_Y$ maps $f(x_0)$ to the identity.
 \item[(5)] The Albanese variety $\Alb(X)$ is the dual abelian variety of the Picard group
 \[\Pic^0(X):=H^1(X,\mathcal O_X)\slash H^1(X,\mathbb Z), \]
 i.e.,
 $\Alb(X)=\Pic^0\bigl(\Pic^0(X)\bigr)$.
 Identifying $\Pic^0(X)$ with the connected component of the identity of the group of line bundles, the functorial homomorphisms of Albanese variety~$f_*$ is the dual of the pullback map.
\end{enumerate}

\begin{remarkk}
 Notice that in the complex setting, $H^1(X,\O_X)$ can be seen as sheaf cohomology or alternatively as the Dolbeault cohomology group of $(0,1)$-forms.
\end{remarkk}



\subsubsection{Algebraic definition}
The interpretation of the Albanese variety as the dual abelian variety of \smash{$\operatorname{Pic}^0_{X\slash S}$}
allows a more algebraic definition.
 Indeed, when $\X\xrightarrow{g} S$ be a smooth, projective, morphism with connected geometric fibers over a scheme $S$, Grothendieck proved that there exists a smooth $S$-scheme, locally of finite presentation, $\Pic_{\X\slash S}$ representing the relative Picard functor and an abelian scheme (in particular, smooth and proper) $\Pic^0_{\X\slash S}$ over $S$ whose fibers $\Pic^0_{\X_s\slash k(s)}$ are the connected components of the identity in $\Pic_{\X_s\slash k(s)}$. See, for example, \cite[Section~9]{fantechi2006fundamental}, or \cite[Chapter~8]{bosch2012neron} and references therein.
Then by \cite[Theorem~3.3]{grothendieck1961technique} (see also \cite{milneAV,mochizuki2012topics}), we can consider the dual abelian scheme
\[\Alb^0_{\X\slash S}:=\Pic^0_{\Pic^0_{\X/S}\slash S}=\bigl(\Pic^0_{\X\slash S}\bigr)^\vee.\]
If $g\colon \X\to S$ has a section $\sigma_S\colon S\to\X$, then we get an Albanese map to $\Alb^0_{\X\slash S}$ constructed in the following way.

Let $\mathcal P$ be the \emph{unique} Poincar\'e line bundle on $\X\times\Pic^0_{\X\slash S}$ which trivializes along $\sigma_S$ in the sense of \cite[Section~4]{bosch2012neron}.

Then we define the Albanese morphism
\[
 a_{\X\slash S}\colon \ \X \longrightarrow \Alb^0_{\X\slash S},\qquad
 \big(U\xrightarrow{f} \X\big) \longmapsto (f\times\mathrm{id})^*\mathcal P.
 \]
Given \smash{$\X\xrightarrow{f} \Y$} a morphism of smooth, proper, geometrically connected $S$-varieties and $\sigma_S\colon S\allowbreak\to \X$ a section, it follows from the universal property of the Albanese scheme that there exists a~unique homomorphism $f_*$ of $S$-abelian schemes making the following diagram commute:
\[
\begin{tikzcd}
\X\ar[r,"f"]\ar[d,"a_{\X\slash S}"'] & \Y\ar[d,"a_{\Y\slash S}"]\\
\Alb^0_{\X\slash S}\ar[r,"f_*"] & \Alb^0_{\Y\slash S},
\end{tikzcd}
\]
where $a_{\X\slash S}$ and $a_{\Y\slash S}$ sends the sections $\sigma_S\colon S\to \X$ respectively $f\circ \sigma_S \colon S\to \Y$ to the identity element. Alternatively, $f_*$ can be described as the dual (in the category of abelian $S$-schemes, see \cite[Section~I.8]{milneAV}) of the pullback
$f^*\colon \Pic^0_{\Y\slash S}\to \Pic^0_{\X\slash S}$.


\subsubsection{Self-duality for smooth curves} As explained, for example, in \cite[Section~9.5.26]{fantechi2006fundamental}, when $\C\to S$ is a family of smooth, projective, geometrically connected genus $g>0$
curves over a Noetherian base, there exists an \emph{self-duality isomorphism}
$\Alb^0_{\C\slash S}\to \Pic^0_{\C\slash S}$.

 If we have a section $\sigma_0\colon S\to \C$, then the self-duality has a very explicit description using the Abel--Jacobi map $A_{\sigma_0}\colon\C\to\Pic^0_{\C/S}$ given by
\[
 A_{\mathcal \sigma_0}\colon \ \C \longrightarrow \Pic^0_{\C\slash S},\qquad
 (T\to\C) \longmapsto \mathcal O_{\C_T}(p_T -\sigma_0\rvert_T),
\]
where $p_T\colon T\to\C\times_S T$ is the induced section of $\C_T=\C\times_S T$.
The self-duality isomorphism is simply given by the pullback along the Abel--Jacobi map. The isomorphism does not depend on the choice, nor on the existence of the section, see, for example, \cite{esteves2002autoduality} for a proof.

Notice that $A_{\sigma_0}$ coincides with the Albanese map defined using the Poincar\'e line bundle on~${\X\times\Pic^0_{\X\slash S}}$ rigidified along the section $\sigma_0$.

\subsection{Refinement of the moduli spaces}

Let $\M=\M_{g,m}\bigl(Y|D^\pm,\beta,\bfw\bigr)$ be the moduli space of logarithmic stable maps to $Y=\mathbb P_X^1(\O\oplus L)$ considered in the previous section and let us now assume that $L\in\Pic^0(X)$: we have $c_1(L)=0$.


\subsubsection[Morphism to Alb(X)]{Morphism to $\boldsymbol{\operatorname{Alb}(X)}$}
Using the tangency orders $\bfw$, we consider the map
\[ a_\bfw\colon \ (x_i)\in X^n\longmapsto\sum_{i=1}^n w_i a_X(x_i)\in\Alb(X). \]
As $\sum w_i=0$, this map does not depend on the choice of a base point in $X$ and is uniquely defined. Indeed, for $z_0$, $z_0'$ two choices of base point and $a_X$, $a'_X$ the corresponding morphisms, we have that $a'_X= t_{a'_X(z_0)}\circ a_X$ and thus
\[
 \sum_{i=1}^n w_i a'_X(x_i)= \sum_{i=1}^n w_i (a_X(x_i) -a'_X(z_0))= \sum_{i=1}^n w_i a_X(x_i).
\]
Now, consider the morphism from the moduli space to the Albanese variety of $X$ by composing with the evaluation map:
\[ \kappa\colon\ \M\xrightarrow{\ev} X^n\xrightarrow{a_\bfw} \Alb(X)\]
defined by
\[%\label{eq:refinement morphism}
\bigl(\C\xrightarrow{f} X,(p_j)_{j=1}^m,(q_i)_{i=1}^n,\O_{\C}(\alpha)\cong^{\lambda} f^*L\bigr)\mapsto (x_i:=f\circ q_i)_{i=1}^n\mapsto \sum_{i=1}^n w_i a_X(x_i).
\]

Recall that $(q_i)_{i=1}^n$ are marking with logarithmic structure while $(p_j)_{j=1}^m$ are standard schematical points on $\C$.

\begin{lem}\label{lem:mapphibeta}
Let $\beta\in H_2(X,\ZZ)$ be the homology class of a complex curve. The bilinear map
\[ \phi \colon\ \alpha\otimes\omega\in H^{0,1}(X)\otimes H^{1,0}(X)\longmapsto \int_\beta \alpha\wedge \omega\in\CC \]
induces a morphism
\[ \varphi_\beta \colon\ \Pic^0(X)\longrightarrow \Alb(X). \]
\end{lem}
\begin{proof}
We use the isomorphisms $H^{0,1}(X)=H^1(X,\O_X)$ and $H^{1,0}(X)=H^0(X,\Omega_X)$ as well as the Hodge decomposition $H^2(X,\CC)=H^{0,1}(X)\oplus H^{1,0}(X)$. Through the above isomorphisms, the map $H^1(X,\ZZ)\to H^1(X,\O_X)$ yielding $\Pic^0(X)$ is actually the composition of $H^1(X,\ZZ)\to H^1(X,\CC)$ followed by the projection onto $H^{0,1}(X)$.
The bilinear map $\phi$ induces a~linear map~${\phi^*\colon H^{0,1}(X)\to H^{1,0}(X)^*}$. To prove that it descends to a map $\Pic^0(X)\to\Alb(X)$, we need to show that
\[ \phi^*\bigl(H^1(X,\ZZ)\bigr)\subset H_1(X,\ZZ)\subset H^0(X,\Omega_X)^*. \]
Let $\omega\in H^{1,0}(X)$ and $\gamma\in H^1(X,\ZZ)$. We write $\gamma\otimes\CC=\gamma^{0,1}+\gamma^{1,0}$ for its Hodge decomposition in $H^2(X,\CC)$. We have
\begin{align*}
\phi^*(\gamma)(\omega) = {}& \int_\beta \gamma^{0,1}\wedge\omega \qquad\text{(by definition)}\\
={} & \int_\beta \bigl(\gamma^{0,1}+\gamma^{1,0}\bigr)\wedge\omega
= \beta\cap (\gamma\cup\omega)
= (\beta\cap\gamma)\cap\omega,
\end{align*}
where the second equality follows from the fact that \smash{$\int_\beta \gamma^{1,0}\wedge\omega=0$}, since $\beta$ is the class of a~complex curve. The last line ensures that $\phi^*(\gamma)(\omega)$ is a period of $\omega$ since $\beta\cap\gamma\in H_1(X,\ZZ)$. Therefore, as desired, $H^1(X,\O_X)\to H^0(X,\Omega_X)^*$ descends to a map
\[ \varphi_\beta \colon\ H^1(X,\O_X)/H^1(X,\ZZ)\longrightarrow H^0(X,\Omega_X)^*/H_1(X,\ZZ). \tag*{\qed}
\]\renewcommand{\qed}{}
\end{proof}

\begin{remarkk}
Giving a non-degenerate positive bilinear map $H^1(X,\O_X)\otimes H^0(X,\Omega_X)\to\CC$ amounts to providing a polarization, i.e., an ample line bundle $\L$, on the Abelian variety $\Pic^0(X)$, see, for example,~\cite{mumford1974abelian}. In an Abelian variety, a polarization then induces a map
\[ A_\L\colon\ L\in\Pic^0(X)\longmapsto t_L^*\L\otimes \L^{-1}\in\Pic^0\bigl(\Pic^0(X)\bigr)=\Alb(X). \]
Assume $X$ is projective with hyperplane class $h\in H^2(X,\ZZ)$ giving a K\"ahler form on $X$. Then we have a polarization given by the non-degenerate positive bilinear form
\[ \phi\colon\ \alpha\otimes\omega\longmapsto \int_X h^{n-1}\wedge\alpha\wedge\omega. \]
Assume the class $\beta$ is obtained by intersecting sufficiently many hyperplane sections: $\beta$ is Poincar\'e dual to $h^{n-1}$. Then, the previously defined morphism~$\varphi_\beta$ is actually the morphism~$A_\L$ provided by the above polarization on $\Pic^0(X)$.
\end{remarkk}
We claim that the morphism $\kappa$ is in fact constant to $\varphi_\beta(L)$. To prove the latter, we start with the case of smooth curve, and then use it to deal with the case of nodal curves.

\begin{prop}\label{prop-image-of-L-is-constant}
The morphism $\kappa\colon\M\rightarrow \operatorname {Alb}^0(X)$ defined above is constant equal to $\varphi_\beta(L)$.
\end{prop}

\begin{proof}
Let us first look at a geometric point $s=\operatorname{Spec}(\mathbb C)\to\M$ of the moduli space such that~the source curve $C$ is smooth. We claim that in such case
$\kappa(s)=f_*f^*L$,
where~${\O_C \bigl(\sum_{i=1}^n w_i q_i\bigr)\cong f^*L}$ (by definition of the moduli functor) and
$f_*\colon\Pic^0(C)\to\operatorname{Alb}^0(X)$ is the morphism obtained dualizing $f^*\colon\Pic^0(X)\to\Pic^0(C)$ via the self-duality isomorphism of~$\Pic^0(C)$.

To see that, choose $p_0$ a point in $C$ and $z_0=f(p_0)$; then the functoriality of the Albanese morphism for compatible choices of base point gives
\[f_*(f^*L)=f_*\biggl(\O_C\biggl(\sum_{i=1}^n w_i q_i\biggr)\biggr)=f_*\biggl(\sum_{i=1}^n w_i A_{p_0}(q_i)\biggr)= \sum_{i=1}^n w_i a_X(x_i)\]
with the last equality following from the commutativity.

Since we are working over $\mathbb C$, we may use that for $X$ (resp.\ $C$), we have
 \[ \Pic^0(X)=H^1(X,\O_X)/H^1(X,\ZZ) \qquad \text{and}\qquad \Alb(X)=H^0(X,\Omega_X)^*/H_1(X,\ZZ). \]
Through Hodge theory, $H^1(X,\O_X)$ is isomorphic to the Dolbeault cohomology group $H^{0,1}(X)$. In particular, by functoriality, the pullback map $f^*\colon\Pic^0(X)\to\Pic^0(C)$ is in fact induced by~${f^*\colon H^1(X,\O_X)\to H^1(C,\O_C)}$, which is the pullback at the level of $(0,1)$-forms. Using that $H^0(X,\Omega_X)=H^{1,0}(X)$, the composition $f_* f^*$ is induced by the $\CC$-linear map
 \[ \varphi_C \colon\ H^{0,1}(X) \xrightarrow{f^*} H^{0,1}(C)\simeq H^{1,0}(C)^*\xrightarrow{f_*}H^{1,0}(X)^*, \]
where the second arrow is the dual to the pullback map for holomorphic forms. The middle isomorphism is provided by the Poincar\'e duality. Finally, if $\alpha\in H^{0,1}(X)$ is a $(0,1)$-form on $X$ and $\omega\in H^{1,0}(X)$ an holomorphic $1$-form, we have
 \[ \varphi_C(\alpha)(\omega) = \int_C f^*\alpha\wedge f^*\omega = \int_C f^*(\alpha\wedge\omega) = \int_{f(C)}\alpha\wedge\omega= \int_\beta \alpha\wedge\omega, \]
where the penultimate equality is push-pull formula.

Now, assume that $C$ is nodal; we write $\nu\colon\bigsqcup C_V\to C$ for the normalization map and $f_V\colon C_V\to X$ for the restriction of $f\circ\nu$ to the component $C_V$ of the normalization. Since each $C_V$ is now smooth, the discussion above tells us that the morphism
\[
f_{V\ast}f_V^*\colon\ \Pic^0(X)\to\Pic^0(C_V)\to\operatorname{Alb}^0(X)
\]
is equal to $\varphi_{\beta_V}$, where $\beta_V=f_{V\ast}[C_V]$. Furthermore, since $\sum \beta_V=\beta$, we have
\begin{align*}
\varphi_{\beta}=\sum \varphi_{\beta_V}\colon\ \Pic^0(X) \to\prod_V\Pic^0(C_V)\to\operatorname{Alb}^0(X),\qquad
 L \to \nu^*f^*L\to \sum_V f_{V\ast}(f_V^*L).
\end{align*}

To conclude the proof, we only need to argue that
\[ \sum_V f_{V\ast}(f_V^*L)=\kappa(s).\]

When $C$ is singular, $f^*L\cong^{\lambda}\O_C(\alpha)$ for $\alpha\in H^0\bigl(C, \bar{M}_C^{\rm gp}\bigr)$ such that
\[\O_C(\alpha)\rvert_{C_V}=\O_{C_V}\biggl(\sum_{e\vdash V}s_{e,V}(\alpha)q_e\biggr),\]
satisfying the balancing condition $\sum_{e\vdash V}s_{e,V}(\alpha)=0$. Here $s_{e,V}(\alpha)$ is the outgoing slope of $\alpha$ at~$V$. Notice that $\alpha$ induces an orientation on the edges of $\Gamma$. We denote with $w_e$ the slope along an edge with this induced orientation. Then if $V$, $W$ are the two vertices adjacent to a~bounded edge $e$, we will have $s_{e,V}(\alpha)=-s_{e,W}(\alpha)$ with one of the two being $w_e$.

 We know from the analysis for the smooth case that
 \[f_{V\ast}(f_V^*L)=\sum_{e\vdash V}s_{e,V}(\alpha) a_X(f_V\circ q_i).\]

 From the observation above and the fact that $f_V(q_e)=f_W(q_e)$ since the maps are induced by the maps on the nodal curves, it follows that when we consider $\sum_Vf_{V\ast}(f_V^*L)$ the contributions coming from the nodes cancel out and we obtain the required identity.
\end{proof}

\begin{expl}
 Assume $X=E$ is an elliptic curve. We have $\Alb(E)\cong\Pic^0(E)\cong E$. For the homology class $\beta=a[E]$, the associated map $\varphi_\beta$ is just the multiplication by $a$.
\end{expl}

\begin{remarkk}
For a family \smash{$C\slash S\xrightarrow{F} X\times S$} with non-smooth family of source curves, it is not obvious how to define
\[F_*\colon\ \Pic^{[0]}_{C\slash S}\to\operatorname{Alb}^0(X)\times S.\]
Indeed, \smash{$\Pic^{[0]}_{C\slash S}\to S$} is not an abelian scheme, as it is no longer proper, and we do not have an Abel--Jacobi map. There are two possible ways to fix the situations.
\begin{itemize}\itemsep=0pt
\item We can consider the compactified Jacobian \smash{$\Pic^{[0]}_{C\slash S}\subseteq \overline{\mathbb J}_{C/S}^{\mathcal E_{\pi},\sigma}$} parametrizing rank one, torsion free simple $(\mathcal E_{\pi},\sigma)$-quasi-stable sheaves on the fibers of $C\to S$ for $\mathcal E_{\pi}$ the canonical degree~$0$ polarization on $C\slash S$ in the sense of \cite[Section~2.1]{melo2015compactifications}. We refer the reader to~\cite{melo2015compactifications} and to~\cite{MeloRapagnettaViviani} for existence and properties of the compactifications.
Then a family version of the argument in \cite[Section~5]{MeloRapagnettaViviani} shows that there is an isomorphism of $S$-group
\[
\Pic^{[0]}_{C\slash S}\to \Pic^0\bigl(\overline{\mathbb J}_{C/S}^{\mathcal E_{\pi},\sigma}\bigr)
\]
and we can then define $F_*$ as the dual of $F^*$ composed with the natural inclusion of \smash{$\Pic^{[0]}_{C\slash S}$} with the compactified Jacobian.

\item Alternatively, one should in future be able to use the theory of logarithmic abelian varieties and of logarithmic Picard group developed by Molcho--Wise \cite{molchowiselog}. They show that~$\operatorname{LogPic}_{C^v\slash S}\to S$ (where the superscript $C^v$ means we are considering $C\to S$ endowed with the vertical log structure)
is a proper group stack over $\operatorname{LogSch}$
whose fibers are logarithmic abelian varieties in the sense of Kajiwara, Kato and Nakayama \cite{kajiwara2008logarithmic}.
It is expected that many classical results about abelian varieties admit generalization for logarithmic abelian varieties. In this spirit, in a forthcoming paper \cite{delignepair} Molcho--Ulirsch--Wise will show that $\operatorname{LogPic}_{C^v\slash S}$ is self dual in the category of logarithmic Abelian variety, admits a~Abel--Jacobi map and is universal for morphism to abelian groups over log schemes.
Once such a theory will be fully developed, considering the composition
 \[
 \Pic^{[0]}_{C\slash S}\hookrightarrow \operatorname{LogPic}_{C^v\slash S}\xrightarrow{F_*}\Alb(X)\times S
 \]
we would obtain the desired extension.
\end{itemize}

Once the extension $F_*$ is defined, one can look at the closed points to give an explicit description of $F_*F^*L$ and verify that this coincide with $\sum_{i=1}^n w_i a_X(f\circ q_i)$.
Since we do not need the extension at the level of universal Jacobian but only the map from $\M$, we do not fill in the technical details left out in this remark.
\end{remarkk}



\subsubsection{Refinement by correlation}\label{sec:correlationrefinement}
%\tnote{Still two moduli spaces ...}
Let us go back to the setting of $\M(X,L)$ where we fix the ramification profile. To simplify the notation, for a given point in $\M$, we set $x_i=f\circ q_i$. We are interested in the case where the contact orders have a non-trivial common divisor, i.e., there is some $1\neq\delta|\gcd(w_i)$; the we can consider the morphism
\[ a_{\bfw/\delta}\colon\ (x_i)\in X^n\longmapsto \sum_{i=1}^n \frac{w_i}{\delta}a_X(x_i). \]
Composing with the evaluation map, we define
\begin{gather}
 \kappa^\delta \colon\ \M \longrightarrow T^\delta(L,\beta)\subset\Alb(X), \nonumber\\
 \bigl(f\colon C\to X,(p_j)_{j=1}^m,(q_i)_{i=1}^n, \O_C(\alpha)\rvert_{[f]}\cong^{\lambda}f^*L\bigr) \longmapsto \sum_{i=1}^n\frac{w_i}{\delta} a_X (x_i)
.\label{eq:maptotorsion}
\end{gather}
In fact, $\kappa^\delta$ composed with the multiplication by $\delta$ is the previously defined $\kappa$. As $\kappa$ is constant equal to $\varphi_\beta(L)$ by Proposition \ref{prop-image-of-L-is-constant}, this ensures that $\kappa^\delta$ has values in $ T^\delta(L,\beta)$, the set of $\delta$-roots of $\varphi_\beta(L)$. The latter is a torsor under the finite group $\Tor_\delta(\Alb(X))$ of $\delta$-torsion elements in~$\Alb(X)$. In the particular case where $L=\mathcal O_X$, $T^\delta(X,\O_X)$ is exactly the set of $\delta$-torsion elements in the Albanese variety $\Alb(X)$.



As $L$ does not possess any canonical $\delta$-root, Proposition \ref{prop-image-of-L-is-constant} no longer applies for $\kappa^\delta$, which has no reason to be constant anymore. The morphism $\kappa^\delta$ defined in \eqref{eq:maptotorsion} thus determines a decomposition of the moduli space of logarithmic maps to $\mathbb P_X(\mathcal O_X\oplus L)$ into components, according to the value taken by $\kappa^\delta$. This refinement is called \textit{correlation}, as it stems from relations between the images of the points inside the Albanese variety. The image $\kappa^\delta(f)$ is called a \textit{correlator} and is usually denoted by $\theta$.

\begin{remarkk}
 The preimages $\bigl(\kappa^\delta\bigr)^{-1}(\theta)$ of elements $\theta\in T^\delta(L,\beta)$ are called components of the moduli space but they are not necessarily connected.
\end{remarkk}

In particular, it follows from the discussion above that we have the following.
\begin{prop}
 Let $\M$ be the moduli space of logarithmic stable maps to $Y=\PP_X(\mathcal O\oplus L)$ with the divisorial log structure defined by $D^- + D^+$. Assume that $L\in\Pic^0(X)$ and that $\delta|\gcd(w_i)$.
 Then we have a splitting of the virtual class
 \[ \vir{\M}=\sum_{\theta\in T^\delta(L,\beta)}\vir{\M^\theta},\]
 where $T^\delta(L,\beta)$ is the $\operatorname{Tor}_{\delta}(\Alb(X))$-torsor of $\delta$-roots of $\varphi_\beta(L)$.
\end{prop}

\begin{defi}\label{def:fullrefined}
We encompass the information of the correlated virtual classes in the \textit{full correlated virtual class}
\[ \fvir{\M}{\delta} = \sum_{\theta\in T^\delta(L,\beta)} \vir{\M^\theta}\cdot (\theta),\]
which is an element in the group algebra $\QQ[\Alb(X)]$ with cycle coefficients and support over the torsor $T^\delta(L,\beta)$.
\end{defi}

\subsubsection{Correlated Gromov--Witten invariants}
The projection $\PP_X(\O\oplus L)\to X$ identifies $D^-$ and $D^+$ with $X$. The evaluation map hence takes the following form:
\[
\ev\colon\ \M_{g,m}\bigl(Y|D^\pm,\beta,\bfw\bigr) \longrightarrow Y^m\times X^n,
 \]
we integrate pullback of cohomology classes over the full correlated virtual class to get \textit{correlated Gromov--Witten invariants}: for $\gamma_1,\dots,\gamma_m\in H^\bullet(Y,\QQ)$ and $\widetilde{\gamma}_1,\dots,\widetilde{\gamma}_n\in H^\bullet(X,\QQ)$, we set
\[ \bigl\langle\bigl\langle\gamma_1,\dots,\gamma_m,\widetilde{\gamma}_1,\dots,\widetilde{\gamma}_n\bigr\rangle\bigr\rangle^\delta_{g,\beta,\bfw} = \int_{\fvir{\M}{\delta}} \prod_1^m\ev_i^*(\gamma_i)\prod_1^n\widetilde{\ev}_i^*(\widetilde{\gamma}_i), \]
which is an element in $\QQ[\Alb(X)]$ with support on the torsor $T^\delta(L,\beta)$. As the computations in the elliptic case from Section \ref{sec-computation-local-invariants} illustrate, distinct correlators may provide different Gromov--Witten invariants, so that the refinement is non-trivial.






\subsection{Deformation invariants and relations of correlated classes}
Let $(\X,\L)\xrightarrow{p} B$ be a family of smooth projective varieties over $B$ and let $\mathcal L\in \Pic^0(\mathcal X\slash B)$ (in fact log smoothness is sufficient).
Let $\Y=\mathbb P(\O\oplus\L)\to B$ the associated family of $\mathbb P^1$-bundles, with the natural boundary structure $\D=\mathcal D^- +\mathcal D^-$.
Fix $\beta$ an effective curve class on $\X$ relative to~$B$ and
fix contact order $\bfw=(w_1,\dots,w_n)$ of $(q_1,\dots,q_n)$ with the boundary.
As in the previous section, the results of \cite{abramovich2014stable,gross2013logarithmic,maulik2023logarithmic} we have that the moduli stack
$\M_B:=\M_{g, m}(\mathcal Y|\D,\beta_b, \bfw)$
parametrizing diagrams of logarithmic maps
\[
\begin{tikzcd}
C\ar[r,"F"]\ar[d]\ar[d] & (\mathcal Y|\mathcal{D})\ar[d]\\
S\ar[r] & B
\end{tikzcd}
\]
is a Deligne--Mumford stack, proper over $B$ end endowed with a perfect obstruction theory $R\pi_*F^*T^{\log}_{\mathcal Y\slash B}$ relative to $B$.

As for the case of absolute stable maps,
for $B$ regular and connected, the basic properties of the virtual class construction \cite{behrend1997intrinsic,manolache2012virtual} ensure that
the degree of the class $\ev^*\gamma\cap\vir{\M_b}$ does not depend on $b\in B$; we also refer the reader to \cite[Appendix~A]{mandel2020descendant} and references therein for more details on the deformation invariance.

As $\M_B$ is a special case of Corollary~\ref{prop:logmapsdegeneration}, we still have the description of the moduli space in terms of logarithmic trivializations. As in the case of $B=\operatorname{Spec}(\mathbb C)$, we have the map
\[\kappa_B\colon\ \M_B\to\Alb(\mathcal X\slash B)\]
defined by
\[\bigl(\C\xrightarrow{f}\X,(p_j)_{j=1}^m,(q_i)_{i=1}^n,f^*\L\cong\O_{\C}(\alpha)\bigr)\rightarrow \sum_{i=1}^n w_i a_{\X\slash B}(f\circ q_i)\]
which is constant with valued $\varphi_{\beta_b}(\L_b)$ on fibers of $B$. In particular, if $\delta$ divides $\gcd(w_i)$, we get a refinement morphism
\[\kappa^\delta_B\colon\ \M_B\longrightarrow T^\delta_B(\X,\L)\subset \Alb^0(\X/B),
\]
where $T^\delta_B(\X,\L)$ is a the torsor under the $B$-group $\operatorname{Tor}_{\delta}\bigl(\Alb^0(\mathcal X\slash B)\bigr)$ of $\delta$-roots of $\varphi_{\beta_b}(\L_b)$ over~$B$. This torsor is not necessarily trivial, but we can always choose $B'\to B$ \'etale such that $T^\delta_B(\X,\L)\times_B B'\cong T^\delta_{B'}(\X',\L')$ trivializes.

Thus, working \'etale locally on $B$ we can always assume that
$T^\delta_B(\X,\L)$ has a section over~$B$, i.e., it is trivial. Then $\M_B$ split into connected component $\M_B^\theta$ indexed by the sections $\theta\colon B \to T^\delta_B(\X,\L)$.

\begin{theo}\label{thm:definvariants}
The correlated Gromov--Witten invariants are deformation invariant, i.e., let~${\X\to B}$ and $\M_B\to B$ as before and assume that $T^\delta_B(\X,\L)$ has a trivializing section $\theta$. Then~for any such section the degree of classes \smash{$\ev^*\gamma\cap\vir{\M_b^{\theta(b)}}$} for $\gamma\in\text{H}^*(\Y^{n+m})$ does not depend on~${b\in B}$.
\end{theo}

\begin{proof}
 Since \smash{$\M_B^\theta\to \M_B$} is open, the perfect obstruction theory \smash{$R\pi_*F^*T^{\log}_{\mathcal Y\slash B}$} relative to $B$ restricts to $\M_B^\theta$. The results then follows once again from the propertied of the virtual class construction \cite{behrend1997intrinsic,manolache2012virtual} as explained for example in \cite[Appendix~A]{mandel2020descendant}.
\end{proof}

By choosing suitable families of deformations of $\Y\to B$, we can exploit Theorem~\ref{thm:definvariants} to find non-trivial identities among the correlated classes.

The morphism $\varphi_\beta\colon\Pic^0(X)\to\Alb(X)$ induces a morphism between their $\delta$-torsion elements. We thus have the subgroup $\varphi_\beta\bigl(\Tor_\delta\bigl(\Pic^0(X)\bigr)\bigr)\subset\Tor_\delta(\Alb(X))$. In particular, the support $T^\delta(L,\beta)$ of the correlated Gromov--Witten invariants is stable by the action of $\varphi_\beta\bigl(\Tor_\delta\bigl(\Pic^0(X)\bigr)\bigr)$. Furthermore, we have the following.

\begin{theo}\label{theo-torsor-invariance}
 The correlated invariants \smash{$\langle\langle\gamma_1,\dots,\gamma_m,\widetilde{\gamma}_1,\dots,\widetilde{\gamma}_n\rangle\rangle^\delta_{g,\beta,\bfw}$} are invariant under the action of $\varphi_\beta\bigl(\Tor_\delta\bigl(\Pic^0(X)\bigr)\bigr)$.
\end{theo}

\begin{proof}
We consider the family over $B=\Pic^0(X)$ induced by the universal line bundle: the fiber over $L\in\Pic^0(X)$ is $\PP(\O\oplus L)$. Thus, in this case, the torsor appearing before Theorem~\ref{thm:definvariants}~is
\[ T^\delta_B(\X,\L) = \big\{ (L,\theta)\in\Pic^0(X)\times\Alb(X) \text{ s.t. } \delta\theta=\varphi_\beta(L)\in\Alb(X)\big\}.
 \]
To get a section, we use the base change $\mu_\delta\colon L\mapsto L^{\otimes\delta}$ from $\Pic^0(X)$ to itself. We thus have
\[ \mu_\delta^*T^\delta_B(\X,\L) = \big\{ (L,\theta)\in\Pic^0(X)\times\Alb(X) \text{ s.t. } \delta\theta=\varphi_\beta\bigl(L^{\otimes\delta}\bigr)\in\Alb(X)\big\}. \]
This torsor now has a section provided by $L\mapsto\theta_0(L)=\varphi_\beta(L)$.
Deformation invariance tells us that for any choice of $R$, the correlators $\theta_0(L)$ and $\theta_0(L\otimes R)$ provide the same invariants. However, the latter are correlators for the a priori distinct $\PP^1$-bundles associated to $L^{\otimes\delta}$ and~${L^{\otimes\delta}\otimes R^{\otimes\delta}}$. Assuming $R$ is of $\delta$-torsion, these $\PP^1$-bundles are the same. Furthermore, the correlators differ by
\[ \theta_0(L\otimes R)-\theta_0(L) = \varphi_\beta(L\otimes R)-\varphi(L)=\varphi_\beta(R), \]
yielding the result.
\end{proof}

Finally, we have the following relations between the classes corresponding to different levels of refinement.

\begin{lem}\label{lem-unrefinement}
The multiplication by $d$ in $\Alb(X)$ induces a morphism of $\QQ[\Alb(X)]$ denoted by~$\m{d}$. The different full refined virtual classes are related as follows:
\[ \m{\frac{\delta}{\delta'}}\bigl(\fvir{\M}{\delta}\bigr) = \fvir{\M}{\delta'}. \]
\end{lem}

\begin{proof}
 The relation merely comes from the fact that multiplication by $\delta/\delta'$ provides a surjection from $\delta$-roots of $\varphi_\beta(L)$ to $\delta'$-roots of $\varphi_\beta(L)$, so that if $\theta'\in T^{\delta'}(L,\beta)$, then we have
 \[ \M^{\theta'}=\bigsqcup_{(\delta/\delta')\theta=\theta'}\M^{\theta}. \tag*{\qed}
 \]\renewcommand{\qed}{}
\end{proof}


\setcounter{section}{4}
% Original source lines 1278--1701
\section{Computation of local invariants}
\label{sec-computation-local-invariants}

\subsection{General considerations and statement}

Our goal in this section is to compute the correlated GW invariants in the case where $X=E$ is an elliptic curve, the genus of the source curve is $g=1$, and with a unique interior point constraint. These \textit{local invariants} will be used in combination with Theorem~\ref{theo-refined-decomposition-surjective-case} to obtain an explicit computation algorithm and derive regularity results for general correlated invariants in Section \ref{sec-regularity}.

We consider $Y=\PP(\O\oplus L)$ for $L\in\Pic^0(E)$; fix an homology class $\beta=a[E]\in H_2(E,\ZZ)$ and a vector $\bfw=(w_1,\dots,w_n)$ of tangency orders. The moduli space
\[ \M(a,\bfw)=\M_{1,1}\bigl(Y|D^\pm,a[E],\bfw\bigr) \]
is the moduli space of log stable maps as in the previous section.
For $\delta|\mathrm{gcd}(w_i)$, we saw that it decomposes into components $\M^\theta(a,\bfw)$.
The correlators $\theta$ satisfy
$
\delta\theta\equiv \varphi_{a[E]}(L)=a\lambda$,
where~\smash{$\lambda=\varphi_{[E]}(L)$} is the image of the line bundle $L\in\Pic^0(E)\simeq E$ through the isomorphism~$\varphi_{[E]}$. We denote the torsor by $T^\delta(L,a)$.

\subsubsection{Uncorrelated case} We recall the computation of the non-refined invariant
\[ \langle \pt_0,1_{w_1},\pt_{w_2},\dots,\pt_{w_n}\rangle_{1,a[E],\bfw}=
\int_{\vir{\M(a,\bfw)}}\ev_0^*(\pt)\prod_2^n\ev_i^*(\pt). \]

\begin{lem}\label{lem-non-refined-computation}
The uncorrelated relative GW invariant has the following value:
\[ \langle \pt_0,1_{w_1},\pt_{w_2},\dots,\pt_{w_n}\rangle_{1,a[E],\bfw}=a^{n-1}\sigma(a)\cdot w_1^2, \]
where $\sigma(a)=\sum_{d|a}d$ is the sum of divisors function.
\end{lem}

\begin{proof}
This elementary formula is computed, for instance, in \cite{blomme2022floor}. Briefly, it may be obtained as follows. Genus $1$ parametrized curves in $Y$ realizing the class $a[E]+b(\bfw)\big[\PP^1\big]$ come from degree~$a$ covering maps $f\colon C\to E$. Up to translation, these are group homomorphisms. We~may fix the translation parameter using the marked point $p_0$, mapped to a fixed point in $E$.
Then we see that these are in bijection with the index $a$ sublattices of $\pi_1(E)\simeq\ZZ^2$, of which there are~$\sigma(a)$.

 Given one of the $\sigma(a)$ covers $f\colon C\to E$, enhancing to a map to $Y=\PP(\O\oplus L)$ amounts to find a section of $f^*L$, with poles and zeros of prescribed order $w_i$ and fixed image in $E$.

We denote by $y_i$ the %position of
marked points in $C$ (with not trivial contact order) and let $x_i=f(y_i)$ their image in $E$. For $2\leqslant i\leqslant n$, there are $a$ possible choices for each $y_i$ given that $x_i$ is fixed. The position of $y_1$ is determined by the relation $\sum w_iy_i\equiv f^*(\lambda)$ in $C$, i.e., $y_1$ is a $w_1$-root of~${f^*(\lambda)-\sum_{i=2}^nw_iy_i}$.
There are $w_1^2$ possible choices for $y_1$, yielding the result.
\end{proof}

\subsubsection{Statement} We now consider the correlated invariants
\[ \langle\langle \pt_0,1_{w_1},\pt_{w_2},\dots,\pt_{w_n}\rangle\rangle^\delta_{1,\beta,\bfw}
= \int_{[\![\M(a,\bfw)]\!]^\delta}\ev_0^*(\pt)\prod_2^n\ev_i^*(\pt), \]
which is an element of the group algebra $\QQ[E]$ with support on $T^\delta(L,a)$.



Before giving a closed formula for the full local correlated invariant, we need to introduce certain functions with values in $\QQ[E]$.

First, we consider the average of torsion elements
\[\vartheta_d:=\frac{1}{d^2}\sum_{d\theta\equiv 0}(\theta) \in\QQ[E].\]
They satisfy \smash{$\vartheta_{d_1}\vartheta_{d_2}=\vartheta_{\mathrm{lcm}(d_1,d_2)}$}.
We then define the following functions: for $d|\delta$,
\[
 \boldsymbol{\vartheta}_\delta(d) = \prod_p (\vartheta_{p^{\nu_p(d)}}-\mathds{1}_{\nu_p(d)<\nu_p(\delta)}\vartheta_{p^{\nu_p(d)+1}} ),
 \]
where the products is over primes and $\nu_p$ is the $p$-adic valuation. The $\mathds{1}$ indicates that the second term for each factor of the product only appears if $\nu_p(d)<\nu_p(\delta)$.

\begin{expl}
 If $d=\delta$, so that $\boldsymbol{\vartheta}_\delta(\delta) =\vartheta_\delta$. If $\delta=p^v$ is the power of a prime number $p$, the values of $\boldsymbol{\vartheta}_{p^v}$ for its divisors $1,p,\dots,p^{v-1},p^v$ are
 $\vartheta_1-\vartheta_p,\vartheta_p-\vartheta_{p^2},\dots, \vartheta_{p^{v-1}}-\vartheta_{p^v}$ and $\vartheta_{p^v}$.
\end{expl}

Next, we consider the variants of the arithmetic function $\sigma(a)=\sum_{d| a}d$ defined as follows: \smash{$\overline{\sigma}^d(a)=\sigma(a/d)$} if $d|a$ and $0$ else. Combining both, we define the following function, with values in $\QQ[\Tor_\delta(E)]$
\[ \boldsymbol{\sigma}_\delta(a)=\sum_{d|\delta}\overline{\sigma}^{\delta/d}(a)\boldsymbol{\vartheta}_\delta(d). \]
For concrete computations, it may be convenient to express $\boldsymbol{\sigma}_\delta(a)$ as a linear combination of the~$\vartheta_d$ for $d|\delta$ but with different coefficients. To do so, consider the following function on $\NN$:
\[
\Upsilon^\delta_d(a) = \prod_p \bigl(\overline{\sigma}^{p^{\nu_p(d)}}-\mathds{1}_{\nu_p(d)<\nu_p(\delta)}\overline{\sigma}^{p^{\nu_p(d)+1}} \bigr)\bigl(p^{\nu_p(a)}\bigr).
\]
In the above product of function, the argument $a$ is factorized over prime numbers, so that for~$a=\prod p^{\nu_p(a)}$ and $d=\prod p^{\nu_p(d)}$,
\[ \overline{\sigma}^d(a)=\prod\overline{\sigma}^{p^{\nu_p(d)}}\bigl(p^{\nu_p(a)}\bigr). \]
The following identity is proven in the proof of Theorem~\ref{theo-expression-local-correlated-invariant}. Thanks to multiplicativity, it suffices to check it for powers of primes, where it stems from a summation by part
\[ \boldsymbol{\sigma}_\delta(a)=\sum_{d|\delta}\Upsilon_d^\delta(a)\vartheta_{\delta/d}. \]
Doing a summation by parts for each prime numbers is called a \textit{multiplicative summation by parts}.

The torsor consists of $\delta$-roots of $a\lambda$. However, $a\lambda$ has no canonical $\delta$-root. The furthest we can naturally define is the $\gcd(a,\delta)$-root \smash{$\frac{a\lambda}{\gcd(a,\delta)}$}. Then, pick $\theta_0$ to be any choice of root such that $\frac{\delta}{\gcd(a,\delta)}\theta_0 = \frac{a}{\gcd(a,\delta)}\lambda$. A correlator $\theta_0$ as before is called a \textit{special correlator}. Theorem~\ref{theo-torsor-invariance} ensures that the result does not depend on the choice of the latter.

One way to construct special correlators is as follows: if $\lambda_0$ satisfies $\delta\lambda_0=\lambda$, one may take~${\theta_0=a\lambda_0}$. Indeed, one has
\[ \frac{\delta}{\gcd(a,\delta)}(a\lambda_0) = \frac{a}{\gcd(a,\delta)}\delta\lambda_0 = \frac{a}{\gcd(a,\delta)}\lambda. \]

\begin{theo}\label{theo-expression-local-correlated-invariant}
The full local correlated invariant has the following expression:
\begin{equation}\label{eq-local-invariant}
\langle\langle \pt_0,1_{w_1},\pt_{w_2},\dots,\pt_{w_n}\rangle\rangle^\delta_{1,a[E],\bfw} = a^{n-1}w_1^2 \boldsymbol{\sigma}_\delta(a)
\cdot (\theta_0),
\end{equation}
where $\theta_0$ is any special correlator, i.e., satisfies \smash{$\frac{\delta}{\gcd(a,\delta)}\theta_0 = \frac{a}{\gcd(a,\delta)}\lambda$}.
\end{theo}

Without assuming Theorem~\ref{theo-expression-local-correlated-invariant}, it can be checked from the definition of $\boldsymbol{\sigma}$ that
\begin{equation}\label{eq-sigma-invariance}
 \boldsymbol{\sigma}_\delta(a) = \boldsymbol{\sigma}_\delta(a)\vartheta_{\delta/\gcd(a,\delta)},
\end{equation}
so that the right-hand side of \eqref{eq-local-invariant} does not depend on which $\theta_0$ we choose. It may thus be replaced by \smash{$\vartheta_{\delta/\gcd(a,\delta)}\cdot(\theta_0)=\dfk\big[\frac{1}{\delta/\gcd(a,\delta)}\big]\bigl(\frac{a}{\gcd(a,\delta)}\lambda\bigr)$}.

Assuming instead Theorem~\ref{theo-expression-local-correlated-invariant}, equation \eqref{eq-sigma-invariance} is in fact a consequence from Theorem~\ref{theo-torsor-invariance}, which tells us that the correlated invariant is invariant when multiplying by an element of \[\varphi_{a[E]}(\Tor_\delta(E))=a\cdot\Tor_\delta(E)=\Tor_{a/\gcd(a,\delta)}(E).\]
See the second part of the proof of Lemma \ref{lem-k=f(tor)} for a proof of the second equality. In particular, it is also invariant by multiplication by $\vartheta_{\delta/\gcd(a,\delta)}$.

With the above formulation, the correlated invariant appears as a refined version of the uncorrelated invariant, and the refinement takes the form of $\boldsymbol{\sigma}_\delta$ replacing $\sigma$.


\subsubsection{Applications} Before going to the proof, we present some applications.

\begin{expl}
 If $a$ is coprime with $\delta$, every $\overline{\sigma}^{\delta/d}(a)$ vanishes except for $d=\delta$, and we thus recover that
$\boldsymbol{\sigma}_\delta(a) = \sigma(a)\vartheta_\delta$.
 In other words, the curves are equally spread among the correlators.
\end{expl}

\begin{expl}
 Let $J_2$ be the second Jordan function, defined by $J_2(p^\alpha)=p^{2\alpha-2}\bigl(p^2-1\bigr)$. It counts the number of elements of order $n$ in $(\ZZ/n\ZZ)^2$.

 The $(0)$-coefficient of $\boldsymbol{\vartheta}_\delta(d)$ is equal to the product of $(0)$-coefficients for each prime $p$. The~latter is equal to $\frac{1}{p^{2\nu_p(d)}}\!-\!\frac{1}{p^{2\nu_p(d)+2}}$ if $\nu_p(d)\!<\!\nu_p(\delta)$ and $\frac{1}{p^{2\nu_p(\delta)}}$ else. This matches the values of~the~function \smash{$\frac{J_2(\delta/d)}{\delta^2}$}, which is also multiplicative. Therefore,
 we get that
 \[ \langle \pt_0,1_{w_1},\pt_{w_2},\dots,\pt_{w_n}\rangle^{\theta_0}_{1,a[E],\bfw} =a^{n-1}\left(\frac{w_1}{\delta}\right)^2\sum_{d|\delta}J_2(d)\overline{\sigma}^d(a). \]
 However, this application is partially a lie since this computation is actually the first step toward proving Theorem~\ref{theo-expression-local-correlated-invariant}.
\end{expl}

From $\m{\delta/\delta'}\bigl(\fvir{\M(a,\bfw)}{\delta}\bigr) = \fvir{\M(a,\bfw)}{\delta'}$, we deduce that the functions $\boldsymbol{\sigma}_\delta$ satisfy
\[ \m{\delta/\delta'}(\boldsymbol{\sigma}_\delta(a)) = \boldsymbol{\sigma}_{\delta'}(a),\]
which may also be checked directly from the definition of $\boldsymbol{\sigma}_\delta$.

Using Theorem~\ref{theo-expression-local-correlated-invariant}, we immediately get the quasi-modularity result for the generating series of correlated invariants. Here, we extend the notion of quasi-modularity for functions with value in a vector space, here chosen to be $\CC[E]$. In the finite-dimensional case, it just means that all coordinate functions are quasi-modular forms. See Section \ref{sec-modularity} for more details.

\begin{coro}
The following generating series is quasi-modular for $\Gamma_0(\delta)$
\[ \sum_a \langle\langle \pt_0,1_{w_1},\pt_{w_2},\dots,\pt_{w_n} \rangle\rangle^\delta_{1,a[E],\bfw}\sfq^a.\]
\end{coro}

\begin{proof}
 The generating series of $\sigma$ is the first Eisenstein series $E_2(\sfq)$, known to be a quasi-modular form. The result thus follows from the quasi-modularity of the generating series $\sum_a \overline{\sigma}^d(a)\sfq^a = E_2(\sfq^d)$ for the congruence subgroup $\Gamma_0(d)$.
\end{proof}

The rest of the section is dedicated to the proof of Theorem~\ref{theo-expression-local-correlated-invariant}, by refining the proof of Lemma \ref{lem-non-refined-computation}. We proceed in several steps. The first is to study the curves coming from a~common covering map $f\colon C\to E$. Summing over covering maps in Section \ref{sec-0-coeff}, we are then able to find a closed expression for the $(\theta_0)$-coefficient. Multiplicativity properties and an induction relation are thus sufficient to prove the formula from Theorem~\ref{theo-expression-local-correlated-invariant}.

 \subsection{Contribution of a fixed cover}

 Let us consider a fixed cover $f\colon C\to E$, which is a group homomorphism choosing the marked point $p_0$ and its image as neutral element. Its kernel $\ker f$ has cardinality $a$, the degree of the covering. Let $f^*\colon E\to C$ be the dual map between the curves seen as their Picard groups. The lifting condition writes itself $\sum_1^n w_jy_j=f^*(\lambda)$, where $\lambda\in E$ corresponds to the line bundle $L$. We have the following group morphism:
 \[
 C^n \longrightarrow E^{n-1}\times C, \qquad
 (k_j) \longmapsto \left(f(k_2),\dots,f(k_n),\sum_1^n w_jk_j \right)
 \]
 with its kernel $K(f)=\big\{ (k_j)\in C\times(\ker f)^{n-1} \text{ s.t. } \sum_1^n w_jk_j=0 \big\}$. In particular, it sits in the following exact sequence, from which we see it has cardinality $w_1^2 a^{n-1}$
 \[ 0\to \Tor_{w_1}(C) \to K(f) \to (\ker f)^{n-1} \to 0. \]

 For the cover $f$, the set of curves matching the constraints $(x_j)$ is in bijection with the following $K(f)$-torsor
 \[ S_{\lambda,x_2,\dots,x_n}=\left\{ (y_j)\in C^n \text{ s.t. }\forall \ 2\leqslant j\leqslant n \ f(y_j)=x_j \text{ and }\sum_1^n w_jy_j=f^*(\lambda)\in C\right\}, \]
also having cardinality $w_1^2a^{n-1}$. We now wish to refine the above description. To do so, we use the correlator function
 \[
 \kappa^\delta \colon\ C^n \longrightarrow E,\qquad
 (y_j) \longmapsto \sum_1^n \frac{w_j}{\delta}f(y_j).
 \]
 The correlators are the elements $\theta\in E$ satisfying $\delta\theta=f_*f^*(\lambda)=a\lambda$. Among them, recall we have the \textit{special correlators} satisfying $\frac{\delta}{\gcd(a,\delta)}\theta_0=\frac{a}{\gcd(a,\delta)}\lambda$ and that $a\lambda_0$, where $\delta\lambda_0=\lambda$ is a~special correlator.

\begin{lem}\label{lem-k=f(tor)}
 The image of $K(f)$ via the correlator function is $\kappa^\delta(K(f)) = f(\Tor_\delta(C))$. Furthermore, it contains \smash{$\Tor_{\delta/\gcd(a,\delta)}(E)$}.
\end{lem}

\begin{proof}
 The equality follows from the definitions. Assume that $(k_j)\in K(f)$. Then we have that $\sum_1^n w_jk_j=0$ and it follows that $\sum_1^n\frac{w_j}{\delta}k_j\in\Tor_\delta(C)$ and consequently $\kappa^\delta((k_j))=f\bigl(\sum_1^n\frac{w_j}{\delta}k_j\bigr)\in f(\Tor_\delta(C))$. Conversely, let $f(s)\in f(\Tor_\delta(C))$ with $s\in\Tor_\delta(C)$. Let $k_1$ be such that $s=\frac{w_1}{\delta}k_1$, which exists because $C$ is divisible, and $k_j=0\in\ker f$ if $j\geqslant 2$. We have
 \[ f(s)=f\left(\frac{w_1}{\delta}k_1\right)=\kappa^\delta(k_1,0,\dots,0)\qquad \text{and}\qquad\sum_1^n w_jk_j=\delta s=0. \]
 Thus, we have the reverse inclusion $f(\Tor_\delta(C)\subset\kappa^\delta(K(f))$.

 We now prove that $\Tor_{\delta/\gcd(a,\delta)}(E)\subset f(\Tor_\delta(C))$. Using the dual morphism $f^*$, we start from $f^*(\Tor_\delta(E))\subset\Tor_\delta(C)$.
 We now apply the morphism $f$ and use that the composition $f\circ f^*\colon E\to E$ is the multiplication by $a$, so that $a\cdot\Tor_\delta(E)\subset f(\Tor_\delta(C))$.
 We now claim that~${a\cdot\Tor_\delta(E)=\Tor_{\delta/\gcd(a,\delta)}(E)}$, which concludes the proof:
 \begin{itemize}\itemsep=0pt %[label=$\triangleright$]
 \item If $x$ is $\delta$-torsion, we have \smash{$\frac{\delta}{\gcd(a,\delta)}ax=\frac{a}{\gcd(a,\delta)}\delta x=0$}, so that we have the inclusion
 \[
 a\cdot \smash{\Tor_\delta(E)\subset\Tor_{\delta/\gcd(a,\delta)}(E)}.
 \]
 \item As $|\Tor_u(E)|=u^2$, the short exact sequence
 \[ 0\to \Tor_a(E)\cap\Tor_\delta(E)=\Tor_{\gcd(a,\delta)}(E) \to \Tor_\delta(E) \xrightarrow{a\cdot} a\cdot\Tor_\delta(E) \to 0, \]
 ensures that they have the same cardinality $\bigl(\frac{\delta}{\gcd(a,\delta)}\bigr)^2$.\hfill $\qed$
 \end{itemize}\renewcommand{\qed}{}
\end{proof}

The following proposition gives a description of the correlators in the image of the torsor~$S_{\lambda,x_2,\dots,x_n}$.

\begin{prop}\label{prop-solution-torsor}
The correlators achieved by the solutions, i.e., the set $\kappa^{\delta}(S_{\lambda,x_2,\dots,x_n})$, form a $f(\Tor_\delta(C))$-torsor which contains the special correlators $\theta_0$. Solutions are uniformely spread among the correlators.
\end{prop}

\begin{proof}
 The solutions form the $K(f)$-torsor $S_{\lambda,x_2,\dots,x_n}$. Applying the correlator function, we immediately get a torsor under $\kappa^\delta(K(f))$, equal to $f(\Tor_\delta(C))$ by Lemma \ref{lem-k=f(tor)}. We also get that the solutions split evenly among the elements.

 We now need to prove that it contains the special correlators. Special correlators form a~torsor under $\Tor_{\delta/\gcd(a,\delta)}(E)$, which lies in $f(\Tor_\delta(C))$ by Lemma \ref{lem-k=f(tor)}. Thus, it suffices to show it contains one of them.
 Pick $y_2,\dots,y_n$ such that $f(y_j)=x_j$ and choose $y_1$ such that~${\sum_1^n\frac{w_j}{\delta}y_j=f^*(\lambda_0)}$, with $\delta\lambda_0=\lambda$; notice that this always exists as $C$ as well is a divisible group. In particular, as $\delta\lambda_0=\lambda$, we have that $\sum_1^n w_jy_j=f^*(\lambda)$ and $(y_j)$ is indeed an element of $S_{\lambda,x_2,\dots,x_n}$. Moreover, we have that
$\kappa^\delta((y_j))=f\circ f^*(\lambda_0)=a\lambda_0$,
 which is one of the special correlators. Therefore, all special correlators belong to the image torsor.
\end{proof}

To finish this section, we provide a first expression of the local correlated invariant as a sum over the covers $f\colon C\to E$. The cover $f$ corresponds to a sublattice $\Lambda\subset\pi_1(E)\simeq\ZZ^2$. We can find a basis $(e_1,e_2)$ such that $\Lambda=\langle ke_1,(a/k)e_2\rangle$ for some unique $k$ such that $k^2|a$. We say that~$\Lambda$, and by extension the associated cover, is of type $(k,a/k)$.

Let $\vartheta(f)\in\ZZ[E]$ be the element with coefficient $1$ for every element in $f(\Tor_\delta(C))$.

\begin{prop}\label{prop-first-expression-local-correlated-inv}
 The correlated invariant admits the following expression:
 \[ \langle\langle\pt_0,1_{w_1},\pt_{w_2},\dots,\pt_{w_n} \rangle\rangle^\delta_{1,a[E],\bfw} = a^{n-1}\left(\frac{w_1}{\delta}\right)^2\sum_{f\colon C\to E} \gcd(k,\delta)\gcd(a/k,\delta) \vartheta(f)\cdot (\theta_0). \]
\end{prop}

\begin{proof}
We know by Proposition \ref{prop-solution-torsor} that the $w_1^2 a^{n-1}$ solutions for a fixed cover $f\colon C\to E$ uniformely spread among the possible correlators, so that we get some integer multiple of the torsor $\vartheta(f)\cdot(\theta_0)$ indexing the possible correlators. To conclude, we merely need to compute the cardinality of $f(\Tor_\delta(C))$.

The basis $(e_1,e_2)$ diagonalizing the lattice inclusion provides real coordinates on $E$ and $C$ such that $f$ has the following expression:
\[
f \colon \ C\simeq (\RR/\ZZ)^2 \longrightarrow E\simeq (\RR/\ZZ)^2, \qquad
 (u,v) \longmapsto (ku,(a/k)v )
 \]
from which we see that \smash{$|f(\Tor_\delta(C))|=\frac{\delta^2}{\gcd(k,\delta)\gcd(a/k,\delta)}$}. Therefore, each correlator carries
\[
a^{n-1}\left(\frac{w_1}{\delta}\right)^2\gcd(k,\delta)\gcd(a/k,\delta)
\]
 solutions. Summing over covers yields the result.
\end{proof}

Our next goal is to find an expression avoiding the summation over the morphisms $f\colon C\to E$, i.e., the index $a$ sublattices of $\pi_1(E)\simeq\ZZ^2$.



\subsection[Computation of the (theta\_0)-coefficient]{Computation of the $\boldsymbol{(\theta_0)}$-coefficient}
\label{sec-0-coeff}

The next step toward the proof of Theorem~\ref{theo-expression-local-correlated-invariant} is to use Proposition \ref{prop-first-expression-local-correlated-inv} to compute the $(\theta_0)$-coefficient, where $\theta_0$ is a special correlator, belonging to the support of all terms in the sum.

To do so, we need to know the number of sublattices of $\pi_1(E)\simeq\ZZ^2$ of a given type, which is precisely the definition of the Dedekind $\psi$-function: it the unique multiplicative function such that for any prime number $p$, $\psi(p^\alpha)=p^{\alpha-1}(p+1)$. There are $\psi(n)$ index $n$ sublattices of $\ZZ^2$ having type $(1,n)$, which are also known as \textit{primitive} sublattices. Merely dividing by $k$, we get that there are $\psi\bigl(\frac{a}{k^2}\bigr)$ sublattices of type $(k,a/k)$.

\begin{expl}
We have $\psi(2)=3$ as there are $3$ sublattices of $\ZZ^2$ of index $2$. If $(e_1,e_2)$ is a~basis of $\ZZ^2$, these lattices are $\langle 2e_1,e_2\rangle$, $\langle e_1,2e_2\rangle$ and $\langle e_1+e_2,e_1-e_2\rangle$.
\end{expl}

\begin{remarkk}
In particular, as there are $\sigma(a)$ index $a$ sublattices, we get that
\[ \sum_{k^2|a}\psi\left(\frac{a}{k^2}\right)=\sigma(a). \]
This may also be checked using the multiplicativity of $\psi$ and $\sigma$ and the fact that the identity is true on powers of prime numbers.
\end{remarkk}

Using Proposition \ref{prop-first-expression-local-correlated-inv}, we get the following expression.

\begin{lem}
Let \smash{$s_\delta(a)=\sum_{k^2|a}\gcd(k,\delta)\gcd(a/k,\delta)\psi\bigl(\frac{a}{k^2}\bigr)$}. For the level $\delta$ refinement, the~$(\theta_0)$-coefficient is equal to $a^{n-1}\bigl(\frac{w}{\delta}\bigr)^2s_\delta(a)$.
\end{lem}

\begin{proof}
We need to compute the $(0)$-coefficient of \smash{$\sum_{f\colon C\to E}\gcd(k,\delta)\gcd(a/k,\delta)\vartheta(f)$}. As each~$\vartheta$ has by definition coefficient $1$ at $(0)$, and as there are $\psi\bigl(\frac{a}{k^2}\bigr)$ lattices of type $(k,a/k)$, we get the result.
\end{proof}

Before getting to the computation of the full local correlated invariant, we transform the above expression to better suit for our purposes, transforming the sum over $k^2|a$ into a sum over~$d|\delta$ using arithmetic functions more convenient than $\psi$.
\begin{itemize}\itemsep=0pt %[label=$\triangleright$]
 \item We already introduced the sum of divisors function $\sigma(a)=\sum_{d|a}d$, whose Dirichlet generating series is $\zeta(s)\zeta(s-1)$.
 \item Given $d\in\NN$, the Dirichlet generating series of $\overline{\sigma}^d(a)=\mathds{1}_{d|a}\sigma(a/d)$ is
\[ \sum_{a=1}^\infty \frac{\overline{\sigma}^d(a)}{a^s} = \sum_{a'=1}^\infty \frac{\sigma(a')}{(da')^s}=\frac{1}{d^s}\zeta(s)\zeta(s-1). \]
 \item The Dirichlet generating series of the second Jordan function $J_2$ is
\[ \sum_{d=1}^\infty \frac{J_2(d)}{d^s}=\frac{\zeta(s-2)}{\zeta(s)}. \]
\end{itemize}

\begin{lem}
The function $s_\delta(a)$ is fully multiplicative in the following sense:
\[ s_\delta(a)=\prod s_{p^{\nu_p(\delta)}}\bigl(p^{\nu_p(a)}\bigr), \]
where products are over the set of prime numbers. Furthermore, we have the following expression:
\[ s_\delta(a)=\sum_{d|\delta} J_2(d)\overline{\sigma}^d(a). \]
\end{lem}

\begin{proof}
We first prove the multiplicativity: assume that we have the decomposition in products of primes $a=\prod p^{\alpha_p}$ and $\delta=\prod p^{\delta_p}$. The sum over $k^2|a$ writes itself as sum over the families~$(i_p)$ with $2i_p\leqslant \alpha_p$ indexed by the set of prime numbers $\P$
\[ \sum_{k^2|a}\gcd(k,\delta)\gcd(a/k,\delta)\psi\bigl(a/k^2\bigr) = \sum_{2i_p\leqslant \alpha_p}\prod \psi\bigl(p^{\alpha_p-2i_p}\bigr)p^{\min(i_p,\delta_p)+\min(\alpha_p-i_p,\delta_p)}, \]
using the multiplicativity of the Dedekind function $\psi$. We can factor the sum as a product and get the sought multiplicativity.

To prove the desired identity, we compute both Dirichlet generating series. The multiplicativity allows to factor the generating series as a product over prime numbers so that
\[ \sum_{a,\delta=1}^\infty \frac{s_\delta(a)}{\delta^sa^t} = \prod_{p\in\P} \biggl(\sum_{d,2i\leqslant\alpha} \psi\bigl(p^{\alpha-2i}\bigr)p^{\min(i,d)+\min(\alpha-i,d)}p^{-sd-t\alpha} \biggr). \]
To finish the computation of the Dirichlet series for $s_\delta(a)$, the inner sum may be rewritten
\[ \sum_{i,j,k}\psi\bigl(p^i\bigr)p^{\min(j,k)+\min(i+j,k)}p^{-ks-2jt-it}, \]
where $p$ is a prime number and $i,j,k\geqslant 0$. Setting aside $i=0$, $\psi\bigl(p^i\bigr)=p^{i-1}(p+1)$ and we can split according to the value of $k$ and compute the geometric sums, yielding some rational function in $p^{-t}$ and $p^{-s}$.

For the second expression, the Dirichlet series is
\begin{align*}
\sum_{\delta,a} \frac{\sum_{d|\delta} J_2(d)\overline{\sigma}^d(a)}{\delta^s a^t} ={} & \sum_{d,k,a}\frac{J_2(d)\overline{\sigma}^d(a)}{d^sk^sa^t} \qquad(\text{where}\quad\delta=dk) \\
= {}& \sum_{d,k} \frac{J_2(d)}{d^s}\frac{1}{k^s}\frac{1}{d^t}\zeta(t)\zeta(t-1)
= \frac{\zeta(s+t-2)}{\zeta(s+t)}\zeta(s)\zeta(t)\zeta(t-1).
\end{align*}
Using that $\zeta(s)=\prod\frac{1}{1-p^{-s}}$, it can also be expressed as the product of values of a rational function at prime numbers. By a technical but elementary computation, we check that both rational functions are actually the same, so that generating functions have equal coefficients, finishing the proof of the identity.
\end{proof}

Notice that if $\gcd(a,\delta)=1$, we have $\overline{\sigma}^d(a)=0$ for every divisor $d$ of $\delta$. Then, we have $s_\delta(a)=\sigma(a)$. This way, the remaining terms may appear as a correction term when $\gcd(a,\delta)\neq 1$.



\begin{expl}
If $p$, $q$ are distinct prime numbers, we have the following values:
 \begin{gather*}
s_p = \sigma+\bigl(p^2-1\bigr)\overline{\sigma}^p,\qquad
s_{p^2}= \sigma+\bigl(p^2-1\bigr)\overline{\sigma}^p+\bigl(p^4-p^2\bigr)\overline{\sigma}^{p^2},\\
s_{pq} = \sigma + \bigl(p^2-1\bigr)\overline{\sigma}^p + \bigl(q^2-1\bigr)\overline{\sigma}^q+\bigl(p^2-1\bigr)\bigl(q^2-1\bigr)\overline{\sigma}^{pq}.
\end{gather*}
\end{expl}


\subsection{Case of other correlators}

To finish the proof of Theorem~\ref{theo-expression-local-correlated-invariant}, we now study the case where $\theta\neq\theta_0$ is another correlator. We proceed in several steps:
 \begin{itemize}\itemsep=0pt %[label=$\circ$]
 \item First prove that the coefficients do not depend on the precise correlator $\theta$ but only on the order of $\theta-\theta_0$ inside $\Tor_\delta(E)$. To do so, we use the deformation invariance for a suitable family where we deform the base curve $E$.
 \item Then, through Lemma \ref{lem-unrefinement}, we prove an induction relation that enables a formal computation of the correlated invariant for $\theta\neq \theta_0$.
 \item Using the induction relation, we show the multiplicativity of the coefficient functions, so that we may restrict to the case of powers of primes, which we compute also using the induction. Combination of both leads to Theorem~\ref{theo-expression-local-correlated-invariant}.
 \end{itemize}

\subsubsection{Dependence on the order} Assumed that $L=\O$ and choose $\theta_0=0$, so that the torsor of correlators is actually $\Tor_\delta(E)$. For $\theta\in\Tor_\delta(E)$, its order $\omega(\theta)$ is the smallest positive integer $n$ such that $n\theta\equiv 0$.

\begin{lem}
 Assuming $L=\O$, the correlated invariant \smash{$\langle \pt_0,1_{w_1},\pt_{w_2},\dots,\pt_{w_n}\rangle_{1,a[E],\bfw}^\theta$} only depends on the order $\omega(\theta)$.
\end{lem}

\begin{proof}
 Write $E$ as some $E_\tau=\CC/\langle1;\tau\rangle$ for some complex number $\tau$ in the Poincar\'e half-plane. For any \smash{$\gamma=\bigl(\begin{smallmatrix}
 a & b \\ c & d \\
 \end{smallmatrix}\bigr)\in {\rm SL}_2(\ZZ)$}, we have an isomorphism between $E_\tau$ and $E_{\gamma\cdot\tau}$ where \smash{$\gamma\cdot\tau=\frac{a\tau+b}{c\tau+d}$}. The isomorphism is actually given by
 \[ \varphi\colon\ (z \modulo 1,\tau) \longmapsto \frac{z}{c\tau+d} \quad \modulo 1,\gamma\cdot\tau. \]
 Choose a path $\tau(t)$ in the Poincar\'e half-plane going from $\tau(0)=\tau$ to $\tau(1)=\gamma\cdot\tau$. The $\delta$-torsion elements in $E_{\tau(t)}$ are deformed continuously along with $\tau(t)$ as they can be written as $u\tau(t)+v$, where $u,v\in\Tor_\delta(\RR/\ZZ)$. The identification $\varphi\colon E_\tau\to E_{\gamma\cdot\tau}$ then induces some monodromy among torsion elements. More precisely, if $u,v\in \Tor_\delta(\RR/\ZZ)$, we have
 \[ \varphi^{-1}\left(u\frac{a\tau+b}{c\tau+d}+v \right)=u(a\tau+b)+v(c\tau+d) = (au+cv)\tau+(bu+dv). \]
 Therefore, we deduce that the values of the correlated invariants for the torsion elements $\theta=u\tau+v$ and $\gamma\cdot\theta=(au+cv)\tau+(bu+dv)$ are the same. As the action of ${\rm SL}_2(\ZZ)$ on the set of torsion elements of the same order in $\Tor_\delta\bigl((\RR/\ZZ)^2\bigr)$ is transitive, we conclude.
\end{proof}

\subsubsection{Coefficients, induction relation and multiplicativity} For $r|\delta$, let $s_\delta[r](a)$ be such that for any $\theta\in\Tor_\delta(E)$,
\[ \langle\pt_0,1_w,\pt_{w_2},\dots,\pt_{w_n}\rangle_{1,a[E],\bfw}^\theta = a^{n-1}\left(\frac{w}{\delta}\right)^2 s_\delta[\omega(\theta)](a). \]
In particular, we have $s_\delta[1]=s_\delta$, corresponding to $\theta=0$. We also set the global function with values in the group algebra \smash{$S_\delta(a) = \sum_{\delta\theta\equiv 0} s_\delta[\omega(\theta)](a)\cdot (\theta)$}, so that
\[ \langle\langle\pt_0,1_w,\pt_{w_2},\dots,\pt_{w_n}\rangle\rangle^\delta_{1,a[E],\bfw} = a^{n-1}\left(\frac{w}{\delta}\right)^2 S_\delta(a). \]
The definition of $S_\delta$ directly comes from the invariants. Theorem~\ref{theo-expression-local-correlated-invariant} consists in proving that $\frac{S_\delta}{\delta^2}=\boldsymbol{\sigma}_\delta$, as defined before Theorem~\ref{theo-expression-local-correlated-invariant}.

\begin{lem}\label{lem-relation-sdelta}
For every pair $\delta'|\delta$, we have the following relation:
\[ \sum_{r|\delta'}J_2(r)s_\delta[r](a) = (\delta')^2s_{\delta/\delta'}(a). \]
\end{lem}

\begin{proof}
We already have $\m{\delta'}\bigl(\fvir{\M(a,\bfw)}{\delta}\bigr)=\fvir{\M(a,\bfw)}{\delta/\delta'}$. Integrating the point constraints and looking at the $(0)$-coefficient yields
\[ a^{n-1}\left(\frac{w}{\delta/\delta'}\right)^2 s_{\delta/\delta'}(a) = \sum_{\delta'\theta\equiv 0}a^{n-1}\left(\frac{w}{\delta}\right)^2s_\delta[\omega(\theta)](a). \]
The sum is over $\theta$ such that $\delta'\theta\equiv 0$. For any $r|\delta'$, there are $J_2(r)$ elements of order exactly $r$, yielding the desired relation.
\end{proof}

These relations form a system which is triangular for the order given by the divisibility. They are thus sufficient to compute all the functions $s_\delta[r]$. We carry out some examples below.

\begin{expl}
If $\delta=p$ is a prime number, we have
$s_p[1]+\bigl(p^2-1\bigr)s_p[p] =p^2\sigma$.
As we already know that $s_p[1]=s_p=\sigma+\bigl(p^2-1\bigr)\overline{\sigma}^p$, we deduce that
\[ \bigl(p^2-1\bigr)\overline{\sigma}^p+\bigl(p^2-1\bigr)s_p[p] = \bigl(p^2-1\bigr)\sigma, \]
and thus
$s_p[p]=\sigma-\overline{\sigma}^p$.
\end{expl}

\begin{expl}
For $\delta=p^2$, we have the following equations:
\begin{gather*}
 s_{p^2}+\bigl(p^2-1\bigr)s_{p^2}[p]+\bigl(p^4-p^2\bigr)s_{p^2}\big[p^2\big] = p^4s_1 =p^4\sigma, \\
 s_{p^2}+\bigl(p^2-1\bigr)s_{p^2}[p] = p^2s_p =p^2\sigma+\bigl(p^4-p^2\bigr)\overline{\sigma}^p, \\
 s_{p^2} = s_{p^2} = \sigma+\bigl(p^2-1\bigr)\overline{\sigma}^p+\bigl(p^4-p^2\bigr)\overline{\sigma}^{p^2},
 \end{gather*}
 which solves for
\begin{gather*}
 s_{p^2} = \sigma+\bigl(p^2-1\bigr)\overline{\sigma}^p+\bigl(p^4-p^2\bigr)\overline{\sigma}^{p^2}, \qquad
 s_{p^2}[p] = \sigma+\bigl(p^2-1\bigr)\overline{\sigma}^p -p^2\overline{\sigma}^{p^2}, \\
 s_{p^2}[p^2] = \sigma-\overline{\sigma}^p.
 \end{gather*}
\end{expl}

To reduce the computation down to the powers of primes, we use the induction relation to prove the multiplicativity of the functions $s_\delta[r](a)$ and $S_\delta(a)$.

\begin{prop}\label{prop-multiplicativity-S}
The function $s_\delta[r](a)$ and $S_\delta(a)$ are fully multiplicative: for pairs of coprime elements $(r_1,r_2)$, $(a_1,a_2)$ and $(\delta_1,\delta_2)$, we have
\[ s_{\delta_1\delta_2}[r_1r_2](a_1a_2) = s_{\delta_1}[r_1](a_1)s_{\delta_2}[r_2](a_2)\qquad
\text{and}\qquad
S_{\delta_1\delta_2}(a_1a_2)=S_{\delta_1}(a_1)S_{\delta_2}(a_2). \]
\end{prop}

\begin{proof}
We show multiplicativity by induction. It is true for $r=1$ by multiplicativity of $s_\delta$. For the induction step, we use the identity from Lemma \ref{lem-relation-sdelta} and the multiplicativity of $s_\delta$,
\begin{gather*}
(\delta'_1)^2(\delta'_2)^2 s_{\delta_1/\delta'_1}(a_1)s_{\delta_2/\delta'_2}(a_2) = (\delta'_1\delta'_2)^2 s_{\delta_1\delta_2/\delta'_1\delta'_2}(a_1a_2), \\
\sum_{\substack{r_1|\delta'_1 \\ r_2|\delta'_2}} J_2(r_1)J_2(r_2)s_{\delta_1}[r_1](a_1)s_{\delta_2}[r_2](a_2) = \sum_{r|\delta'_1\delta'_2} J_2(r)s_{\delta_1\delta_2}[r](a_1a_2) \\
 \phantom{\sum_{\substack{r_1|\delta'_1 \\ r_2|\delta'_2}} J_2(r_1)J_2(r_2)s_{\delta_1}[r_1](a_1)s_{\delta_2}[r_2](a_2) }{}=\sum_{\substack{r_1|\delta'_1 \\ r_2|\delta'_2}} J_2(r_1)J_2(r_2)s_{\delta_1\delta_2}[r_1r_2](a_1a_2).
\end{gather*}
To get the last sum, we use that each divisor of $\delta'_1\delta'_2$ can uniquely been written as a product of a divisor of $\delta'_1$ and a divisor of $\delta'_2$. If we assume multiplicativity for $r$ a strict divisor of $\delta'_1\delta'_2$, all the terms in the sum except the last one are equal, and we are left with the multiplicativity for~${r=\delta'_1\delta'_2}$, finishing the induction. The multiplicativity of $S_\delta$ follows from the multiplicativity of~$s_\delta[r]$,
\begin{gather*}
\biggl(\sum_{\delta_1\theta_1\equiv 0}s_{\delta_1}[\omega(\theta_1)](a_1)\cdot(\theta_1) \biggr)\biggl(\sum_{\delta_2\theta_2\equiv 0}s_{\delta_2}[\omega(\theta_2)](a_2)\cdot(\theta_2) \biggr) \\
\qquad= \sum_{\delta_1\theta_1\equiv\delta_2\theta_2\equiv 0}s_{\delta_1}[\omega(\theta_1)](a_1)s_{\delta_2}[\omega(\theta_2)](a_2)\cdot(\theta_1+\theta_2).
\end{gather*}
As each $\delta_1\delta_2$ torsion element can be written uniquely as the sum of a $\delta_1$-torsion and a $\delta_2$-torsion element, we conclude.
\end{proof}

\subsubsection{Computation for powers of primes} Because of multiplicativity, we only need to compute the function when $a,r,\delta$ are the power of a common prime $p$. Recall that we defined the elements of $\QQ[E]$: $\vartheta_d=\frac{1}{d^2}\sum_{d\theta\equiv 0}(\theta)$. We also momentarily set \smash{$\vartheta_d^{\mathrm{prim}}=\sum_{\omega(\theta)=d}(\theta)$}.


\begin{prop}\label{prop-expression-powers-primes}
 The function $(\delta,r,a)\mapsto s_\delta[r](a)$ has the following values over powers of primes, for $r>0$
 \[ s_{p^d}[p^r]=s_{p^{d-r}}-p^{2d-2r}\overline{\sigma}^{p^{d-r+1}} = \sum_0^{d-r} J_2\bigl(p^j\bigr)\overline{\sigma}^{p^j} -p^{2d-2r}\overline{\sigma}^{p^{d-r+1}} . \]
 Consequently, we find the following expressions for $S_{p^d}$:
 \[ \frac{S_{p^d}}{p^{2d}} = \sum_{r=0}^{d-1} \bigl(\overline{\sigma}^{p^r}-\overline{\sigma}^{p^{r+1}}\bigr)\vartheta_{p^{d-r}} + \overline{\sigma}^{p^d}\vartheta_1
 = \sigma\vartheta_{p^d} + \sum_1^d \overline{\sigma}^{p^r} (\vartheta_{p^{d-r}}-\vartheta_{p^{d-r+1}}). \]

\end{prop}

\begin{proof}
 By Lemma \ref{lem-relation-sdelta}, we have the following relations:
 \[ p^{2r}s_{p^{d-r}} = \sum_{j=0}^{r} J_2\bigl(p^j\bigr) s_{p^d}\big[p^j\big]. \]
 For $r>0$, making the difference between relations for $r$ and $r-1$, we get
 \[ p^{2r}s_{p^{d-r}}-p^{2r-2}s_{p^{d-r+1}} = J_2(p^r)s_{p^d}[p^r] = p^{2r-2}\bigl(p^2-1\bigr)s_{p^d}[p^r]. \]
 Therefore, after dividing by $p^{2r-2}$, using that \smash{$s_{p^{d-r+1}}=s_{p^{d-r}}+J_2\bigl(p^{d-r+1}\bigr)\overline{\sigma}^{p^{d-r+1}}$}, we finally get
 \[ \bigl(p^2-1\bigr)s_{p^{d-r}} -J_2\bigl(p^{d-r+1}\bigr)\overline{\sigma}^{p^{d-r+1}} = \bigl(p^2-1\bigr)s_{p^d}[p^r]. \]
 Dividing by $p^2-1$ yields the desired expression. Then, we deduce that $S_{p^d}$ has the following expression:
 \begin{align*}
 S_{p^d} ={} & s_{p^d}\vartheta_1 + \sum_{r=1}^d \bigl(s_{p^{d-r}}-p^{2(d-r)}\overline{\sigma}^{p^{d-r+1}} \bigr)\cdot \vartheta_{p^r}^\mathrm{prim} \\
 ={} & s_{p^d}\vartheta_1 + \sum_{r=1}^d \bigl(s_{p^{d-r}}-p^{2(d-r)}\overline{\sigma}^{p^{d-r+1}} \bigr)\cdot \bigl(p^{2r}\vartheta_{p^r}-p^{2r-2}\vartheta_{p^{r-1}} \bigr).
 \end{align*}
 We can now split the sum performing a summation by part to get
 \begin{align*}
 S_{p^d}= {}& \bigl(s_{p^d}-s_{p^{d-1}}+p^{2(d-1)}\overline{\sigma}^{p^d}\bigr)\vartheta_1 +\bigl(\sigma-\overline{\sigma}^p\bigr)p^{2d}\vartheta_{p^d} \\
 & + \sum_1^{d-1} p^{2r}\vartheta_{p^r} \bigl(s_{p^{d-r}} -p^{2(d-r)}\overline{\sigma}^{p^{d-r+1}}-s_{p^{d-r-1}}+p^{2(d-r-1)}\overline{\sigma}^{p^{d-r}} \bigr) \\
 ={} & p^{2d}\overline{\sigma}^{p^d}\vartheta_1 +p^{2d}\sum_1^d \bigl(\overline{\sigma}^{p^{d-r}} - \overline{\sigma}^{p^{d-r+1}} \bigr)\vartheta_{p^r}.
 \end{align*}
 This yields the first expression. Performing a second summation by part yields the second expression.
\end{proof}

We have two expressions of $S_{p^d}$, depending on which feature we wish to emphasize: the $\vartheta_d$ or the $\overline{\sigma}^d$. The first ones are useful to compute in the group algebra since they are projectors, while the second are multiplicative functions that provide quasi-modularity properties.

\begin{proof}[Proof of Theorem~\ref{theo-expression-local-correlated-invariant}]
Combining the multiplicativity from Proposition \ref{prop-multiplicativity-S} and the expression for powers of primes from Proposition \ref{prop-expression-powers-primes} yields and explicit expression for $S_\delta$ and we finally get that $\frac{S_\delta}{\delta^2}=\boldsymbol{\sigma}_\delta$.
\end{proof}
\begin{thebibliography}{99}
\bibitem[1]{abramovich2014stable}
Abramovich D., Chen Q., Stable logarithmic maps to {D}eligne--{F}altings
 pairs~{II},
 \href{https://doi.org/10.4310/AJM.2014.v18.n3.a5}{\textit{Asian~J. Math.}}
 \textbf{18} (2014), 465--488,
 \href{http://arxiv.org/abs/1102.4531}{arXiv:1102.4531}.

\bibitem[2]{abramovich2014comparison}
Abramovich D., Marcus S., Wise J., Comparison theorems for {G}romov--{W}itten
 invariants of smooth pairs and of degenerations,
 \href{https://doi.org/10.5802/aif.2892}{\textit{Ann. Inst. Fourier
 (Grenoble)}} \textbf{64} (2014), 1611--1667,
 \href{http://arxiv.org/abs/1207.2085}{arXiv:1207.2085}.

\bibitem[4]{behrend1997intrinsic}
Behrend K., Fantechi B., The intrinsic normal cone,
 \href{https://doi.org/10.1007/s002220050136}{\textit{Invent. Math.}}
 \textbf{128} (1997), 45--88,
 \href{http://arxiv.org/abs/alg-geom/9601010}{arXiv:alg-geom/9601010}.

\bibitem[5]{birkenhake}
Birkenhake C., Lange H., Complex abelian varieties, 2nd ed., \textit{Grundlehren Math.
 Wiss.}, Vol.~302,
 \href{https://doi.org/10.1007/978-3-662-06307-1}{Springer}, Berlin, 2004.

\bibitem[7]{blomme2022floor}
Blomme T., Floor diagrams and enumerative invariants of line bundles over an
 elliptic curve,
 \href{https://doi.org/10.1112/s0010437x23007285}{\textit{Compos. Math.}}
 \textbf{159} (2023), 1741--1790,
 \href{http://arxiv.org/abs/2112.05439}{arXiv:2112.05439}.

\bibitem[9]{BorneVistoli}
Borne N., Vistoli A., Parabolic sheaves on logarithmic schemes,
 \href{https://doi.org/10.1016/j.aim.2012.06.015}{\textit{Adv. Math.}}
 \textbf{231} (2012), 1327--1363,
 \href{http://arxiv.org/abs/1001.0466}{arXiv:1001.0466}.

\bibitem[10]{bosch2012neron}
Bosch S., L\"utkebohmert W., Raynaud M., N\'eron models, \textit{Ergeb. Math.
 Grenzgeb.}, Vol.~21,
 \href{https://doi.org/10.1007/978-3-642-51438-8}{Springer}, Berlin, 1990.

\bibitem[14]{CN}
Carocci F., Nabijou N., Rubber tori in the boundary of expanded stable maps,
 \href{https://doi.org/10.1112/jlms.12874}{\textit{J.~Lond. Math. Soc.}}
 \textbf{109} (2024), e12874, 36~pages,
 \href{http://arxiv.org/abs/2109.07512}{arXiv:2109.07512}.

\bibitem[16]{esteves2002autoduality}
Esteves E., Gagn\'e M., Kleiman S., Autoduality of the compactified {J}acobian,
 \href{https://doi.org/10.1112/S002461070100309X}{\textit{J.~London Math.
 Soc.}} \textbf{65} (2002), 591--610,
 \href{http://arxiv.org/abs/math.AG/9911071}{arXiv:math.AG/9911071}.

\bibitem[17]{fantechi2006fundamental}
Fantechi B., G\"ottsche L., Illusie L., Kleiman S.L., Nitsure N., Vistoli A.,
 Fundamental algebraic geometry: {G}rothendieck's {FGA} explained,
 \textit{Math. Surveys Monogr.}, Vol. 123,
 \href{https://doi.org/10.1090/surv/123}{American Mathematical Society},
 Providence, RI, 2005.

\bibitem[19]{gross2013logarithmic}
Gross M., Siebert B., Logarithmic {G}romov--{W}itten invariants,
 \href{https://doi.org/10.1090/S0894-0347-2012-00757-7}{\textit{J.~Amer. Math.
 Soc.}} \textbf{26} (2013), 451--510,
 \href{http://arxiv.org/abs/1102.4322}{arXiv:1102.4322}.

\bibitem[20]{grothendieck1961technique}
Grothendieck A., Technique de descente et th\'eor\`emes d'existence en
 g\'eom\'etrie alg\'ebrique. {VI}.~{L}es sch\'emas de {P}icard: propri\'et\'es
 g\'en\'erales, in S\'eminaire {B}ourbaki, {V}ol.~7, Soci\'et\'e
 Math\'ematique de France, Paris, 1995, Exp.~No.~236, 221--243.

\bibitem[22]{kajiwara2008logarithmic}
Kajiwara T., Kato K., Nakayama C., Logarithmic abelian varieties,
 \href{https://doi.org/10.1017/S002776300000951X}{\textit{Nagoya Math.~J.}}
 \textbf{189} (2008), 63--138.

\bibitem[23]{KatoF}
Kato F., Log smooth deformation and moduli of log smooth curves,
 \href{https://doi.org/10.1142/S0129167X0000012X}{\textit{Internat.~J. Math.}}
 \textbf{11} (2000), 215--232.

\bibitem[24]{kim2010logarithmic}
Kim B., Logarithmic stable maps, in New {D}evelopments in {A}lgebraic
 {G}eometry, {I}ntegrable {S}ystems and {M}irror {S}ymmetry ({RIMS}, {K}yoto,
 2008), \textit{Adv. Stud. Pure Math.}, Vol.~59,
 \href{https://doi.org/10.2969/aspm/05910167}{Mathematical Society of Japan},
 Tokyo, 2010, 167--200,
 \href{http://arxiv.org/abs/0807.3611}{arXiv:0807.3611}.

\bibitem[26]{li2001stable}
Li J., Stable morphisms to singular schemes and relative stable morphisms,
 \href{https://doi.org/10.4310/jdg/1090348132}{\textit{J.~Differential Geom.}}
 \textbf{57} (2001), 509--578.

\bibitem[28]{mandel2020descendant}
Mandel T., Ruddat H., Descendant log {G}romov--{W}itten invariants for toric
 varieties and tropical curves,
 \href{https://doi.org/10.1090/tran/7936}{\textit{Trans. Amer. Math. Soc.}}
 \textbf{373} (2020), 1109--1152,
 \href{http://arxiv.org/abs/1612.02402}{arXiv:1612.02402}.

\bibitem[29]{manolache2012virtual}
Manolache C., Virtual pull-backs,
 \href{https://doi.org/10.1090/S1056-3911-2011-00606-1}{\textit{J.~Algebraic
 Geom.}} \textbf{21} (2012), 201--245,
 \href{http://arxiv.org/abs/0805.2065}{arXiv:0805.2065}.

\bibitem[30]{marcuswiselog}
Marcus S., Wise J., Logarithmic compactification of the {A}bel--{J}acobi
 section, \href{https://doi.org/10.1112/plms.12365}{\textit{Proc. Lond. Math.
 Soc.}} \textbf{121} (2020), 1207--1250,
 \href{http://arxiv.org/abs/1708.04471}{arXiv:1708.04471}.

\bibitem[32]{maulik2023logarithmic}
Maulik D., Ranganathan D., Logarithmic enumerative geometry for curves and
 sheaves, \href{https://doi.org/10.4310/cjm.250319031722}{\textit{Camb.~J.
 Math.}} \textbf{13} (2025), 51--172,
 \href{http://arxiv.org/abs/2311.14150}{arXiv:2311.14150}.

\bibitem[33]{melo2015compactifications}
Melo M., Compactifications of the universal {J}acobian over curves with marked
 points, \href{http://arxiv.org/abs/1509.06177}{arXiv:1509.06177}.

\bibitem[34]{MeloRapagnettaViviani}
Melo M., Rapagnetta A., Viviani F., Fourier--{M}ukai and autoduality for
 compactified {J}acobians.~{I},
 \href{https://doi.org/10.1515/crelle-2017-0009}{\textit{J.~Reine Angew.
 Math.}} \textbf{755} (2019), 1--65,
 \href{http://arxiv.org/abs/1207.7233}{arXiv:1207.7233}.

\bibitem[35]{milneAV}
Milne J.S., Abelian varieties~(v2.00), 2008, {a}vailable at
 \url{https://www.jmilne.org/math/}.

\bibitem[36]{mochizuki2012topics}
Mochizuki S., Topics in absolute anabelian geometry~{I}: generalities,
 \textit{J.~Math. Sci. Univ. Tokyo} \textbf{19} (2012), 139--242.

\bibitem[37]{delignepair}
Molcho S., Ulirsch M., Wise J., The logarithmic {D}eligne pairing, {P}rivate
 communication.

\bibitem[38]{molchowiselog}
Molcho S., Wise J., The logarithmic {P}icard group and its tropicalization,
 \href{https://doi.org/10.1112/s0010437x22007527}{\textit{Compos. Math.}}
 \textbf{158} (2022), 1477--1562,
 \href{http://arxiv.org/abs/1807.11364}{arXiv:1807.11364}.

\bibitem[39]{mumford1974abelian}
Mumford D., Ramanujam C.P., Manin J.I., Abelian varieties, \textit{Stud.
 Math.}, Vol.~5, Oxford University Press, 1974.

\bibitem[41]{ogus}
Ogus A., Lectures on logarithmic algebraic geometry, \textit{Cambridge Stud.
 Adv. Math.}, Vol.~178, \href{https://doi.org/10.1017/9781316941614}{Cambridge
 University Press}, Cambridge, 2018.

\bibitem[42]{ranganathan2024logarithmic}
Ranganathan D., Kumaran A.U., Logarithmic {G}romov--{W}itten theory and double
 ramification cycles,
 \href{https://doi.org/10.1515/crelle-2023-0100}{\textit{J.~Reine Angew.
 Math.}} \textbf{809} (2024), 1--40,
 \href{http://arxiv.org/abs/2212.11171}{arXiv:2212.11171}.

\bibitem[43]{RSPW2}
Ranganathan D., Santos-Parker K., Wise J., Moduli of stable maps in genus one
 and logarithmic geometry,~{II},
 \href{https://doi.org/10.2140/ant.2019.13.1765}{\textit{Algebra Number
 Theory}} \textbf{13} (2019), 1765--1805,
 \href{http://arxiv.org/abs/1709.00490}{arXiv:1709.00490}.

\bibitem[44]{ranganathan2020rational}
Ranganathan D., Wise J., Rational curves in the logarithmic multiplicative
 group, \href{https://doi.org/10.1090/proc/14749}{\textit{Proc. Amer. Math.
 Soc.}} \textbf{148} (2020), 103--110,
 \href{http://arxiv.org/abs/1901.08489}{arXiv:1901.08489}.

\end{thebibliography}
\end{document}
